\documentclass[a4paper,11pt,openany]{ltjsbook}
\usepackage{luatexja}
\usepackage{luatexja-fontspec}
\usepackage{amsmath,amssymb,amsfonts,amsthm}
\usepackage{geometry}
\geometry{top=24mm,bottom=24mm,left=24mm,right=24mm,footskip=10mm}
\usepackage{tcolorbox}
\tcbuselibrary{skins,breakable}
\usepackage{hyperref}
\usepackage{fancyhdr}
\usepackage{graphicx}
\usepackage{tikz}
\usetikzlibrary{arrows.meta,calc,positioning,shapes.geometric}

\definecolor{primaryblue}{RGB}{20, 60, 110}
\definecolor{darkslate}{RGB}{35, 45, 60}
\definecolor{accentcrimson}{RGB}{160, 30, 45}
\definecolor{softblue}{RGB}{238, 244, 252}

\hypersetup{
    colorlinks=true,
    linkcolor=primaryblue,
    citecolor=darkslate,
    urlcolor=primaryblue,
    pdftitle={インド・アラビア数字の歴史 —— ゼロの発見と近世記数法の革新},
    pdfauthor={デイヴィッド・ユージーン・スミス、ルイ・チャールズ・カーピンスキー / 教育情報分析研究会}
}

\pagestyle{fancy}
\fancyhf{}
\fancyhead[LE,RO]{\thepage}
\fancyhead[RE]{\leftmark}
\fancyhead[LO]{\rightmark}
\setlength{\headheight}{17.5pt}
\renewcommand{\headrulewidth}{0.4pt}

\title{
    {\huge\bfseries\gtfamily インド・アラビア数字の歴史}\\[14pt]
    {\Large ゼロの発見と近世記数法の革新}\\[8pt]
    {\large —— 原著完全翻訳・数学校閲版 ——}\\[16pt]
    {\normalsize 原題：\textit{The Hindu-Arabic Numerals} (Project Gutenberg eBook No. 22599)}
}
\author{
    \textbf{デイヴィッド・ユージーン・スミス}（David Eugene Smith）\\
    \textbf{ルイ・チャールズ・カーピンスキー}（Louis Charles Karpinski） 著\\[12pt]
    \textbf{教育情報分析研究会} 訳・数理歴史校閲
}
\date{2026年10月 完全決定版}

\begin{document}
\maketitle
\frontmatter
\tableofcontents
\clearpage

\chapter*{解題：『インド・アラビア数字の歴史』の学術的意義}
\addcontentsline{toc}{chapter}{解題：『インド・アラビア数字の歴史』の学術的意義}

本書は、アメリカ数学史研究のパイオニアであるコロンビア大学教授デイヴィッド・ユージーン・スミス（David Eugene Smith, 1860--1944）とミシガン大学教授ルイ・チャールズ・カーピンスキー（Louis Charles Karpinski, 1878--1956）による名著『\textit{The Hindu-Arabic Numerals}』（1911年刊）の完全日本語翻訳・数学校閲版である。

今日、人類のあらゆる科学技術、経済取引、日常的計算の根幹を支えている「$0, 1, 2, 3, 4, 5, 6, 7, 8, 9$」という10個の記号と、それらに位置による重み付けを与える「十進位取り記数法（place value system）」は、古代インドで発明され、中世アラビア世界を経由して西欧中世・ルネサンス期へと伝播した。

本書は、古代インドのアショーカ王碑文（紀元前3世紀）やブラーフミー文字、ナーナー・ガート石刻、グワリオール碑文（876年の最古の確定ゼロ刻印）にまで遡り、アル＝フワーリズミーやアル＝キンディーらイスラム黄金期の数学者たちの貢献、さらにはシルヴェステル2世（教皇ゲルベルト）のアバクス（算盤）革命、レオナルド・フィボナッチの『算盤の書（Liber Abaci）』によるヨーロッパへの本格的受容、そして15世紀の印刷術普及に至る1500年以上にわたるドラマを、膨大な一次史料と碑文・写本証拠に基づいて緻密に論証した記念碑的労作である。

高校までの数学的基礎知識（位取りの概念や初等代数・算術操作）をもつ読者であれば、人類最大の知的発明の一つである「記数法革命」の全貌を鮮明に追体験することができるであろう。

\clearpage
\mainmatter
\chapter{第1章：インド・アラビア数字の起源に関する初期の諸説}
\section*{EARLY IDEAS OF THEIR ORIGIN}

第一章

彼らの起源に関する初期の考え

日常生活で使用される一般的な数値は、比較的最近のものであることが長く認識されてきた。紀元前の記号表記システムの数は、書かれた言語の数とほぼ同じであり、一部の言語には複数のシステムが存在した。例えば、エジプト人は三つの書き方を持ち、それぞれに数値表記を用いていた；ギリシャ人は二つの明確な数値体系を持っており、ローマの数字は一世纪ごとに多少変化していた。今日でも、一般的な数値概念を表現する方法の数は、その科目について学ぶ前に想像するよりも遥かに多いことが分かっており、私たちが使用している一般的な数値が普遍的であるという考えは誤りである。したがって、「アラビア数字」と呼ばれる数値は、キリスト教以前に使用されていた多くのシステムの一つと考えるべきだ。その当時のシステムは他の多くのものと比べても優れておらず、最近の発明であり、ゼロを含まず、使いにくく且つ頻繁には用いられていなかったし、特に将来性も見られていなかった。数世紀後の後になってからこそ、このシステムはビジネスや科学界で地位を得た；そしてギリシャ人が現在の位値概念とゼロが必要なものを開発していたら、今日使っているような数字ではなく、ギリシャ数字を使っていたかもしれない。

世界が最初に使用した数値形態についてここで述べる場所ではない。それらは興味深いものではあるが、科学的な価値はほとんどなかった。ヨーロッパでは、記号表記の発明は一般に地中海東岸で行われたとされ、一世纪前後の重要な時期までさかのぼる。時としてユダヤ人、時としてエジプト人、しかしもっと頻繁には早期の商取引をしていたフェニキア人によるとされる。

当時の学術書籍において、我々が用いる一般的な数字符号がアラビア語の起源であるという考えは古くないものではない。中世とルネサンス期の著者はこれらをインド語のものとして一般に認識し、多くの者は明確にそのインドの起源を述べていた 。他には、カルデア人やユダヤ人がこれらの数字符号をおそらく発明したとする議論があり、右から左へと価値が増すという点から、ローマ語やギリシャ語のシステムもまた同様に適用されるだろうという主張があった。実際には、多くのこれらの著者は記号表記の一般概念を指摘し、イギリスで最初の算術教本作者であるロバート・レコード（約1542年頃）はその言葉から明らかであると述べている。「この芸術を最初に発明したカルデア人は、これらの数字を他のすべての文字と同じように配置した。彼らはあなたが逆順と言っているように書き、また読み方も同じである。これはヘブライ語、カルデア語、アラビア語のすべての書籍において明らかである……一方で、ギリシャ人、ラテン人、そしてヨーロッパのすべての民族は左手から右手へと書き、読みる」。

他には、イタリアのタルタリアやドイツのコベルのような有力な著者が数字符号のアラビア語起源を主張し、また他の者たちはこの問題について判断を保留したり、「野蛮なもの」として片付けたりした。もちろん、アラブ人は自らが発明者であることを主張せず、常にインド人に対する数字符号の形態と位置価値の特徴的な性質に感謝していた。その中でも特に重要な著者は、バグダード黄金時代の偉大な師匠であり、東洋と西洋の数学的古典を最初のアラブ学者として集め、保存し、最終的に欧州の啓蒙期へと伝えた者である。この人物はモハメド・ムサの息子、またはアラブ風に表記するとモハメド・イbn・ムサ・アル・クロズミであり、彼は学問が深く、世界が現在の代数と算術に関する知識を得るのに大きく貢献した人物である。彼についてしばしば話題となることがあるし、彼が著した算術においてアデルハール・オブ・バース（約1130年頃）が翻訳または解釈を試みたものの中で明確に述べているように、数字符号はインド人によって発明されたとされている。これは現在まで続くアラブの後期の著者たちにおいても明白に主張され続けている。実際には、「hindī」（インド語）という用語は彼らにとって算術を指し、単独で「hindī」という形容詞としても使用されている。

アブ・アルヒーン・アルビルーニ（973-1048）は、アラビアの旅行者で学術者であり、『インドに関する大作』『古代歴法論』を著し、「暗号書」も執筆したが、不幸にして失われてしまった。彼はまた、幾多の他の著作物の作者でもあった。アルビルーニは、アラビア語、ペルシャ語、サンスクリット語、ヘブライ語、シリア語に通じており、天文、歴法、数学にも精通していた。『インドに関する大作』においては、その国の言語や習慣について詳細な情報を提供し、当時のヒンドゥー人はアルファベットの文字を使用して数値を表すことをしなかったと明記している。また、数字記号「アンカ」がインド各地で異なる形状を呈していたことも述べており、これはアルファベットも同様であった。『古代諸民族の歴法』においては、幾何級数の和を示し、誤差を避けるためには三つの異なる表現方法を使用することが必要であることを説明している：インドの符号、六十進記数法、後述するアルファベットシステム。さらに、「179,876,755をインド暗号で表す」と述べており、これもまたヒンドゥー語源とみなされる形態である。

アルビルーニ以前には、第十世紀の另一位アラビア語作家モタハッラー・イbn・タヒールが存在し、『創造と歴史の書』を著しており、インド語（ナガリ文字）で表記した多くの数をインドの人々は世界の持続時間として主張しました。フアルトは確信していますが、モタハッラーの時代には現在のアラビア語の数字がまだ使用されていなかったこと、そして学者たちには知られていたものの普及していないインドの文字で表記された数字が現れたことに驚くことはありませんでした。この場合、著者も読者もその数に何か不思議な点を見出すことがなかったでしょう。

アルマズドィ（885?-956）は、広く旅行を重ねた学生で、バグダードからペルシャ、インド、シーロンまで旅し、また中国海を越えてマダガスカル、シリア、パレスチナにも訪れた。彼はあらゆる入手可能な情報源に注意を払ったようであり、ペルシア人、ヒンドゥー人、ローマ人の歴史も検討した。カリフ時代の時期については、『金の草原』という著作から非常に興味深い情報が得られる。彼は、王によって招集されたインドの賢者たちが『シンヒンド』を編纂したと述べている。さらに、歴史家ムハンマド・イブン・アリ・アブディによると、アルマンサールの命令により多くの科学および占星術に関する著作がアラビア語に翻訳され、特に『シンヒンド』（シッダントァ）が翻訳されたと述べている。この名前の意味や拼写については意見が分かれる。コルブロックは最初、『シッダントァ』と『シンヒンド』の関連性を指摘し、「回転する時代」という意味だと説明している。セディロットは、これらの名称を集めた一方で、一般的にヒンドゥー人に関連付けられる科学がギリシャ語起源である可能性があると主張した。カシリは『学問の歴史』（タルイク・アル・フカム）を引用し、この著作は『シンドゥム・インデム』であり、「永遠不変」という意味だと述べている。この古いアラビア語著作におけるアルコワルズミへの参照は興味深いものである。

この『シンヒンド』は、マズド＝イ所说，すべてのインド人が知る天体、星占術、算術およびその他の学問について記述した書籍である。彼はまたアル＝カフワラジミとハバシュを翻訳者として挙げている。アル＝ビルーニは、サインドからカルイファム＝アル＝マンスールへの政治使節団の一員としてバグダードに来たインド人が提供した著作からの二つの他の翻訳についても言及している。その年はヒジャリ154年（西暦771年）である。

何らかの意味での完全な作品としては、987年に書かれた『フィヒルスト』が最初であり、イブン＝アビ・ヤーアクブ＝アル＝ナディムによって著された。これはアラビア文化の歴史にとって基礎的な重要性を持つ。この著作の十章の中で、第七章は本稿において注意を引くものである。その第二節は数学と天文学者について扱っているからである。

最初に言及されるアラビアの作家はアル＝キンディ（800-870年）であり、彼は五冊の算術に関する書籍と四冊のインド数法使用に関する書籍を著した。イスラム教に改宗したユダヤ人のセネド・イブン・アリも、インド数法についての著作の作者として挙げられている。しかし、セネド・イブン・アリに帰属される作品の一部が実際にはアル＝カフワラジミのものである可能性がある。彼の名はその直後に示されているからである。しかしながら、この点について注意すべきは、カシリも同じ作家を算術に関する最も有名な著作の著者として述べていることである。

アル-[S.]＝ウが986年没した際には、同様の学術著作を著しており、他の作家による類似の論文も記述されている。従って、中世において新数字符号のインド源流についての共通認識があったことは明らかである。

ピサのレオナルドについては、その数字符号の導入に関する章で詳しく述べるが、1202年に『 Liber Abbaci』を著した。この著作ではしばしば「九つのインド数字」 を参照しており、中世における共通認識である新数字符号のインド源流を再確認している。

また、我々の言語で最も古い算術に関する文書にも興味が持たれる。アルゴリズムに関する最初の論文の一つは、1240年頃のミノレート僧アレクサンドラ＝デ＝ヴィラ＝ディーによって書かれた『アルゴリズモスの歌』 の解説である。その文章の最初の数行は以下の通りである。

「アルゴリズムとは思慮深き芸術を指すもので、
  タリブは二倍の五つのインド数字を楽しむ。

この本はアルゴリズムまたはオグリムと呼ばれ、俗語として用いられた。そしてこの本は数の技術について論じており、その技術はアルゴリズムとも呼ばれる。インドの王アルゴルがこの芸術を作り出した……アルゴリズムでは我々は十つのインド数字を使用する。」

*       *       *       *       *

\clearpage

\chapter{第2章：位取り記法を伴わない初期のインド数字}
\section*{EARLY HINDU FORMS WITH NO PLACE VALUE}

第2章

初期のインド数字（位取り記法を伴わない）

科学的な天文学の発展において、キリスト教時代の初めに至るまでインドにおける学術進歩は、ギリシャや中国の源流に基づいていたと広く認められていますが、彼らの古い文献は高度な文明を証明し、哲学や文学においても著しい進展があったことを示しています。ギリシャの黄金時代以前から長い間、インド人は思想をリズミカルな形式で表現することが好まれてきました。この最初の詩歌（これは名にふさわしく、また形而上学的な観点から見ても価値がある）はヴェーダであり、2000前紀から1400前紀までのヴェーダ期における賛美 hymn と礼拝の詩歌の集積です。その後に続くブラフマニカ文学は、部分的に儀式的（ブラフマナ[=a]ス）であり、部分的に哲学的（ウパニシャッド）です。私たちの特に興味があるのは、儀式と典礼規則の詮略化された表現であるシュトラスであり、これは祭壇建設に関連する幾何学的な材料を含み、また平方和がまた平方となる多数の有理数の例も含まれています。これらはピタゴラスが生きた前から「ピタゴラスの数」として知られています。ウィンチレーはヴェーダ文学を含むヴェーダ、ブラフマナ[=a]ス、シュトラスを1500前紀から800前紀までと位置づけ、ブエルクがシュトラスに現れるピタゴラスの定理に関する知識は紀元前8世紀まで遡ると主張していることに合意しています。

インドの数字符号が一般的に知られるアラビア数字とされる最も古くからの形式がインドにその源をもつことを確実に断定することはできません。その形はエジプト、東部ペルシャ、中国、または美索不達ミヤ平原で使用されたものから Inspiration を受ける可能性があります。これらの初期の段階については、我々にはほとんど情報がありませんが、インドでの発展に関しては、位置価の基本的な特徴が出現した時期、アラブ文明への導入、そして西へと広がった時期について、多少の確実な情報を有しています。したがって、シンハンドランドにおける数字符号の発生を考えるときには、インド自身以外の多くの可能性のある源を考慮に入れなければならず、これらは最初に数字符号が顕著に現れた以前から存在していました。

古代インドの歴史を検討する者は、信頼できる資料の乏しさに強い印象を受けざるを得ない。人々は自らの種族以外の人々にはほとんど同情を示さず、総合的な文学が全く欠けているため、科学的進歩の全面的な見解は異国の者たちの観察からしか得られない。最初期の時代においても基本学校が存在していたという証拠があるし、七十二の認知された学術分野の中で書記と算術が最も重視されていたことが確認できる。

吠陀期、すなわち紀元前2000年から1400年頃までにおいては、バビロニア、中国、エジプトの早期文明における天文学への同じような関心があった。吠陀自体がその事実を証明している。このような科学的進歩は計算に相当の知識が必要とすることを示唆するが、その方法については全く不明であり、おそらく永遠に解明されないだろう。

仏教の聖典『ラリタヴィスタラ』によれば、菩提薩埵が結婚年齢となったとき、ゴパの父は五百人の求婚者に対する検査を要求し、その科目には算術、書記、琵琶、射撃が含まれていた。彼は他の競争相手をすべて打ち破り、大数学者アーラジュナと対決し、100コティを超える数を表すことを求められた。彼の回答として、$10{53}$まである数名前体系を提示し、さらに$10{421}$までは進めることができると述べた。これはアキメデスの体系を想起させ、西洋が古代数学において東洋にどれほど影響を受けたのかという未解決の問題を示唆している。

エドWIN・オルダール卿は『アジアの光』においてこの部分を言及していないが、仏陀が学問的なヴィシュミトラから訓練を受けたことについては述べている。

\&\&
  于是曰：「已矣，
 数字を数えよ。私に従い、数字の数列を唱え、十万まで達す。
 一、二、三、四、十、百、千。」その後、子供は
 十位、百位、千位と名付け、中断せず、円満な十万に達し、静かに続けた。
 次にコ＝オティ、ナヒュット、ニンナヒュー、カムブド、グンドィカ、ウトパラ、プンダリカ、最後にパドゥマが来る。
 パドゥマはハスタギリの最微塵の砂数を数える方法である。
 しかしそれ以上には別の数字がある：
 カ＝サ、夜空の星を数えるために使用される；
 コ＝オティ＝カ＝サ、海洋の滴を数えるために使用される；
 インガ、円形物の計算法；
 サルヴァニッキチェパ、ガンガの砂すべてを扱うための方法。
 そして最後にアントハ＝カルパスで、単位は十億個のガンガの砂である。より広範な尺度を求めれば、アサンカヤにより算術が昇り、
 十万年の間に全ての世界に降る雨滴の数を含むものとなる。
 それからマハ＝カルパスへと至り、神々は過去と未来を計算するために使用する。

そこでヴィサミトラ＝アチャーヤはその仕事について賛成し、線の長さの「測量」まで聞くことを求めた。この指示を受け、仏陀は次のように追加した：

... 「そして先生！ご容赦いただければ、
 一ヨジャナの端から端までの太陽粒子数を唱えますか？」
そこで、小王は瞬時にその原子数を正確に述べた。
しかしヴィサミトラはその顔を拝し、膝を屈め、「汝よ、師の師を教える者なり－
 あなたこそ、 guru なり！」と叫んだ。

この記述は歴史とは程遠いものである。しかし、仏教時代の数列について伝える古い『ラリタヴィスタラ』の内容を優雅なリズムで表現している。

この前キリスト教以前の時代には、『ヴェダanga』も含まれる。これは「ヴェーダの支えとなる部分」であり、伝統によって伝えられるヒンドゥー文学の一大分野『スミティ』（記憶）の一部である。これらのうち第六は『ジヨティシャ』と呼ばれ、天文に関する短い論文で、紀元前300年頃から書かれたもので、その時代における数値計算の範囲を我々に教えてくれる。

また、ヒンドゥーも十八の古い『シダーダンタ』や天文関連の著作について語り、これらは大半が失われているものの、この証拠を支持している。

しかし、真実の歴史に関しては、ムハンマド時代以前（622年）のインドには存在しない。我々が知る限りの早期文明については、『マハーバラタ』と『ラマヤナ』から得られる情報、硬貨、そして数少ない碑文のみである。

次に述べる未完の資料に基づき、ヒンドゥー・アラビア数字の早期歴史を探求せざるを得ない。第四紀末以前にはインドへの文字導入がなかったことと、その時期以前は言語のみで文献が存在したことから、数値工作は原始民族共通のものであり、sticksやcowries、pebblesを地面に並べる、砂まぶしの板に線を引く、または縄を結ぶといった具象的な方法であったことが疑われる。

このようにして、ヒンドゥー初期の数字は三つの主要なグループに分類される。すなわち、(1) Kharoṣṭhī、(2) Brāhmī、および (3) 言葉と文字による表現形態である。これらを順次考察する。

カルホス＝ティ数字符号は、過去にバクトラ語、インド・バクトラ語、アーライアン語として知られていた铭文や、紀元前4世紀から3世紀にかけて古代ガンドハラ地方、現在の東アフガニスタンと北パンジャブで見られる铭文中に見つかる。その言語のアルファベットは、紀元前4世紀から紀元3世紀までの铭文の中で見られ、右から左へと书かれることからセミト系起源であると考えられている。しかし、これらの最初の铭文中では数字符号は见到されておらず、数词は多くの古代民族と同じように言葉で表されていた可能性がある。そして、紀元前3世紀に强き王アソカの时代になると初めて、発见された铭文の中で数字符号が见られるようになる；しかし、その形式は原始的なマークであり、エジプトやギリシャ、ローマ、あるいは世界の他の地域で見られるようなものである。このアソカ铭文は全部で三十枚ほどあり、インドの広範囲に散在し、しばしば柱に见つかる。これらの铭文は人々が知っていた各种の方言で书かれている。その中にはカルホス＝ティ文字で书かれたものが二枚あり、残りはブラーミ文字のどの形式でもある。カルホス＝ティ铭文中では仅かに四つの数字符号しか见られず、これらはそれぞれ1、2、4、5を表す垂直なマークである：

|     |
|     |
||     |
||||

また、可能として紀元前1世紀のサカ铭文では、より発展した形式で多くの数字符号が见られ、右から左へと书かれるシステムとともに、四、十、二十という三個の计数スケールの迹が见られる。この时期的の数字符号は以下の通りである：

[図表]

このシステムを研究する際には、いくつかの重要な点に注意を払うべきである。

まず第一に、ナナガット形式よりもおそらく早いものではないが、ナナガットの彫刻自体はアーソカ期より後である。四はこのシステムにおいてローマ数字のXに類似しており、おそらく職人が五を表すために三つのマークを削除したか、あるいは五本の棒に一本を横たえたものと考えられる。十については説明が不十分であり、カハルシティ文字のAと類似しているが、なぜこの文字が選ばれたのかは不明である。二十は二つの十の組み合わせであることが明白で、これにより基底のような概念が提示され、九十はおそらくフランス語の quatre-vingt-dix に類似した書き方で表されたものと考えられる。百については説明がないが、ブラヒミ文字の $ta$ や $tra$ の右側に1が付いている点と類似している。二百は百の象形文字の変種であり、二本の垂直線を有する。

このシステムはキリスト教初期のネバテア数字と多くの類似性がある。四は十字で表され、カハルシティ形式は二十を表すために使用されている。さらに、二十のスケールを使用した類似的な方法も見られる。100を表す象形文字は異なるが、他の百の形成法は同じである。この対応関係は単なる偶然ではないと考えるのが適切であるように思われる。

因此、カーヒスト文字からは、私たちが使用する数字の起源を期待することはできません。そこで、インドの第二種類の文字、ブラフマキャラクターに注目します。ブラーマと関連付けられるアルファベットは、インドで知られている複数の中でも最も古いものであり、歴史的な初期から使用されていました。その起源に関する諸説がありますが、これらはまだ広く受け入れられていません。しかし、問題は時間の経過とともに解決する可能性があることを示しています。数字はアルファベットよりも古くないか、少なくともアーソカ王の勅令が刻まれた以前には、文字に刻まれた痕跡はまだ見つかっていません。これらの勅令の一部はナガ・ガット山頂の洞窟内のブラフマキャラクターとカーヒスト文字で刻まれていました。既述の通り、古い著者は数を言葉で記した可能性があり、シヤンの最初期パリ文書でも似たようなことが推測されます。

以下に示す数字は、アーソカ勅令の中で知られている唯一のものであり：

[図版]

これらの紀元前3世紀の断片は、粗末で不十分なものではありますが、私たちが現在使用するシステムの確実な初期形態です。それらは紀元前2世紀に中央インドのプーナから約75マイル離れたナガ・ガット山頂の洞窟内の刻まれた文字で再び見られます。これらの刻まれた文字は南インドの早期アンドラ王朝の記念碑である可能性がありますが、それらの主な興味は含まれる数字にあります。

洞窟はカリアナからジュナルへの山登りの旅人を休憩させるために作られました。その山は道路に位置しています。洞窟内部には入口対面の壁に王['サタヴァハナ]の後裔によって彫刻された家族の肖像が見られます。文字は磨灭していますが、名前はまだ読み取れます。掘削は王ヴィデシリの命令により行われたと思われます。「この記述には、数回のヤグナや宗教的な儀式を行う際の贈り物の一覧が含まれています」という記録があり、30か所以上に数字が見られます。

これらの記述で本当に見られる数字についての議論は多いです。解読が難しく、不確かな部分がありますが、擦過からコピーした以下のものは数値の形かもしれません：

[図版]

この記録自体は、私たち自身に関連する最も古く重要なヒンドゥー語の数値システムを含んでおり、その重要性から24ページに部分的に再現することを許す価値があります。

次に非常に注目すべき数字に関する証拠があり、これは1世紀または2世紀頃まで遡る洞窟記述で完全な形で見られます。ナシクの洞窟記述では、以下のような形式が見られます：

[図版]

その後から、デカルト数値システムが最初の九文字を採用し、ブラーミ文字表記の残りをゼロの追加によって置き換えるまでの間、これらの形態の進行は明確です。したがって、最も知られている標本を概観的に示すことが適切であり、25ページのテーブルでそのようにしています。

[表：インドにおける数値形式の進展]

これらの数字符号について、まず指摘すべきは、表にはゼロが含まれていないことである。実際、引用されたすべての例においてもゼロは存在しなかった。そのため、位置価値を有するような符号は不可能であり、二十や三十、他の十の倍数、百の倍数などはそれぞれ独立した符号が必要であった。これらの数字符号は文字で表記される場合を除き、別々の符号が要求された。

古代インド人は少なくとも20種類のこのような符号を持っていたが、その後その数は著しく増加した。以下に示すのは、百から千までの特定の数を表す方法の例である：

174を表す数字符号
191を表す数字符号
269を表す数字符号
252を表す数字符号
400を表す数字符号
356を表す数字符号

これらの符号は、表にあるものと同様の形態である以下の一桁以下の数を表す数字符号も含められる：

90を表す数字符号
70を表す数字符号

これまで述べてきたカルシュタとブラフマ字体は、ゼロの発明によってシステムが完成化される過程と密接に関連しているため、その考察は次章に移すことが適切である。特に、これらの形態はアラビア語として知られる形態の起源とは直接関係がない。

このように早期の形式を検討した後では、その発生についての問題に焦点を当てるべきである。最初の三種については疑いの余地がない。単一の縦線または横線は単なる一筆画であり、計算者が下ろす棒である。二つの縦線または横線はそれぞれ二筆画または二本の棒を表し、同様に三つの縦線または横線もそのように解釈される。原始的な二つの縦線や横線から、エジプト、ローマ、古代ギリシャ、そして他の諸文明の「2」が発展した。これらの数字符号は以下の形でエジプトの三種類の文字体系に現れる：

象形文字 [2 縦線]
ヘリアティック     [ヘリアティック 2]
デモティック      [デモティック 2]

最後のものは単なる草書体であり、アラビア語 [アラビア 2] のように直角に傾けても2と同様である。原始的な二つの横線から中国の  符号が発展し、150 B.C.から700年頃までのインドで頻繁に見られる符号とほとんど同一である。この草書体は [2 横線を接続] となり、ドイツでは18世紀まで二を表すためにしばしば使用された。これは現代の形式2へと進化し、同様に三つの横線も3となった。

これらの最初の三つの数字符号に関しては、興味深い推測の余地が存する。インド最早期の形式は垂直であった。ナナガット銘文ではこれらは縦に並んでいる。しかしアーソカ铭文やナナガット銘文よりも遥か以前には、中国は最初の三つの数字符号を横書きしていたが、四つ目については縦書きを使用していた。現在、これらの形式はどこから中国に伝わったのか？インドからは絶対にないだろう。彼女自身の遺跡と文献が示すところによれば、インド人がそれらを知る以前にすでに存在していたのである。伝統によると、中国はチベット北部を経由してモンゴルから文明を持ち込んだとされる。原始的な居住地はメソポタミアまたはトルキスタンの oasis である可能性がある。

ではメソポタミアはどのような数字符号を使用していたのか？バビロニア体系は、全体的には単純だが、多くの詳細において非常に複雑であり、現在はよく知られている。特に、一、二、三は垂直の矢印で表わされていた。なぜならば、中国がこれらの数字符号を横書きしたのか？この問題は新たな興味を惹く。実際、これらのバビロニア形式はこの地域の原始的なものではなく、早期sumerian形式は横に並んでいたのである。

これらの事実に対する解釈はいかなるものとすべきであろうか？一つの民族が「一」を縦に書き、別の民族が横に書いたという偶然性のみを主張するべきであるか。これもまた可能であるが、シナのモンゴル侵攻まで遡る部落移動がユーフラテス川流域から始まったり、あるいはターキスタンの共通源から始まったりした場合、その原始的な数字符号を故郷から東へ運んだ可能性も否定できない。人々全体として知っていたり必要とされていた最初の三つの数字だけが移動に伴って伝えられたのであろう。

また、これらの横書きの形態は原始的な棒書きを反映している可能性もある。計算板の前に置かれた棒の最も自然な位置が水平であることを考えると、これがその理由かもしれない。しかし、楔形文字が発展した段階で縦書きの方が容易に書けることが証明された場合、西への移動が始まる頃には既に使用されており、それらはエジプト、ギリシャ、ローマ、地中海諸国の直立した数字の形態やアーソカ時代のインドの文字を生み出したのであろう。

アーソカ以降、あるいは前世紀の商人たちによって、これらの横書きの形態が中国からインドに伝わり、ナナ・ガット洞窟や後の銘文の数字の形態となった可能性もある。これは推測の域を出ないが、更なる彫刻学的研究によりこの仮説が確認される可能性は否定できないのであろう。

以下の翻訳です：

位取り記法を伴わない初期のインド数字について、三桁以上の数には多くの推測が存在する。デモティック文字の一はサンスクリットの一に類似し、二（逆向き）はアラビアの一に類似し、四はナシク洞窟の数字と多少類似性がある、五（逆向き）はカトラパ銭の数字と類似性があり、九は Kushana の碑文の数字と類似性があるなど、類似点が多数存在した。また、ヒエログラフィック文字の五と七はインドの碑文に類似しているとの推測もあったため、これら数字が起源をエジプトに持つという主張者もいた。

一方で、Kharoshthi 数字はかつてバクトラー、インド・バクトラー、アーライーンと呼ばれていた。Cunningham はこれらの数字が東ペルシャのバクトラー文明の文字から派生した可能性があると初めて提案し、彼の説を学術的に支持したのは Sir E. Clive Bayley であり、その後 Canon Taylor が続いた。しかし、類似性は説得力に欠けるものであり、Bayley の図版は彼の理論に基づいており批判を受けた。以下はその仮説の一端である：

| 数字 | ヒンドゥー語 | バクトラー語 | サンスクリット語 |
|------|------------|-------------|----------------|
| 4    | [シンボル]  | [シンボル] = ch | chatur, 拉丁語 quattuor |
| 5    | [シンボル]  | [シンボル] = p  | pancha, ギリシャ語 [p/ente] |
| 6    | [シンボル]  | [シンボル] = s  | [s.]a[s.] |
| 7    | [シンボル]  | [シンボル] = [s.] | sapta |
| （s と [s.]は西北インドで時々逆転している） |

Buehler はこの仮説を否定し、四（4）、六（6）、七（7）、十（10）の四つの場合においては、事実上この推測が支持されないと述べている。

古代バクトラの形式との関連性については、一般的に疑問視されており、多くの数字符号がブラーマ文字に類似していることは広く認められています。自然には、これらはアルファベットの頭文字であると期待されるところがあります。しかし、大部分の数値名詞の古語発音を知る限りでは、この予想とは異なることが明らかになります。次に、これらの符号が古代アルファベットの文字順序を表しているという仮説に回帰します。しかしながら、その文字順序についての知識からは、この推測にも基礎がないことが分かります。したがって、これらの数字符号がアルファベット的であるかどうかは確実ではありませんし、それがアルファベット的である場合でも、現在では選択基準に関する証拠はありません。後の形式については、特定の音節を表すとされるアクラサラ（akṣaras）の固定された数値を持つサスティク語において、これらの符号がアルファベット的な表現であった可能性があるかもしれませんが、これは他の仮説と同じく確実ではありません。ベイリーも、いくつかの形式がフェニキア語的であると考えており、特に二十を円で表す使用はその一例ですが、その類似性は一般的には説得力がないと指摘しています。

また、一部の初期のインド数字符号に中国語の影響がある可能性もあります。

ギリシャ語起源説については、現在の符号がギリシャアルファベットの最初の九文字から派生したとする主張により支持されているとされますが、これはいくつかの文字を曲げたり、切り取ったり、追加したりして理論に適合させるという困難な作業です。この不自然な説明は、ダスイポディウス（コンラード・ラウフス）によって最初に提唱され、後にはヒュエトによって詳細化されました。

キルケラーは、1040年頃のアラビア語（）の初期資料に基づく奇妙な導出を『数秘術』において述べている。彼はアブ・ラジャールの著作から引用し、アラビア語とそのラテン語訳（）を提示し、通常のアラビア文字は円弧の部分から派生したと主張しているという。（）

すべてこれらの対立する理論から、そして西欧の数字と東アジアの数字との間にある実在または想像される類似性から、現在の証拠に基づいて我々はどのような結論を下すべきか？おそらく満足できるものはない。実際、手元に存在する証拠によれば、インドの数字が本質的にインド土着であるという点では何らかの形で全てが一致しているように思われる。なぜこれが事実ではないのか？639年（唐天寶7年）に拉萨を建国した王松贊干布が新アルファベットを制定することができたとすれば、タイ人、僧伽羅人、ビルマ人その他の東アジアの人々が自らの文字を作り出すことも不可能ではない。なぜなら、数字もまた寺院学校や国王、商人組合によって作られた可能性がある。例として、表36に示す東アジアのいくつかの体系には、少数の類似性は見られるものの、それぞれの創始者は他の体系に存在しない独自の形を模索しようと努めたことが明白である。したがって、この証拠に対する公正な解釈としては、人間の心が絶対的に新しい単純な形式を作り出すことは容易ではなく、それが形成するものは他者の心によって形成されたものや耳で聞いたり目に見えて知ったものを自然に模倣することになる。円は世界共通の図形であり、フェニキアとインドで20を意味し、バビロニアでは一度十を意味したとしても驚くべきことではない。したがって、外的な起源を見つけることは不可能であるという十分な理由から、そのようなものは存在しないのである。

本章について総括すると、我々が使用する数字符号の起源に関する絶対的な無意味な説明は多く書かれている。しかし、歴史的証拠を支えずに推測することは、現行の数字符号の逐次的発展について真剣な議論には適さない。

以下に特定の東方系統を示す表を附す。
シアム [図版：数字符号]
ビルマ [図版：数字符号]
マラバール [図版：数字符号]
チベット [図版：数字符号]
セイロン [図版：数字符号]
マラヤラム [図版：数字符号]

ある特定の初期のインド数字符号がどのような理由から使用されたかについては、誰も確かな知識を持っていない。一部の起源は明らかであるが、他の部分の起源はおそらく永遠に明らかにならないだろう。これらの数字符号が何かしらの祭司や教師、または職業団体によって発明され、王の命令により、あるいはある寺院の神秘主義の一環として発明された可能性がある。無論のところ、その起源はいかなるものであれ、これらは他の多くの古代システムと同様に優れておらず、ゼロを用いて記述される現在の中国システムと同等であり、ゼロが導入されていない場合でも西へ進展する機会はほとんどなかった。この相対的に遅く導入された符号なしのゼロがなければ、これらの数字符号の勝利的な進展はあり得なかった。ゼロの文字をインドの起源とする証拠はこれほど強力なものはないだろうし、第四章においてさらに詳細に論じる。

*       *       *       *       *

この証拠は、ゼロの文字がインドの起源であることを強く示しているため、第四章で更なる参照を行う。

\clearpage

\chapter{第3章：位取り記法を伴う後期のインド数字}
\section*{LATER HINDU FORMS, WITH A PLACE VALUE}

第三章

位取り記法を伴う後期インド数字

ゼロと位取り記法を有する完全なヒンディー語の数字符号について述べる前に、第19ページで言及された第三の体系、すなわち単語や文字による形式について考慮しなければならない。位取り記法を持つ単語の使用は、キリスト教 Era 6世紀頃から始まったと考えられている。多くの天文学および数学の手引き書において、またしばしば他の著作で日付を挙げる際に、数は特定の物体や概念の名前によって表される。例えば、「空」（śūnya）または「天空間」（āmbara ākāsa）がゼロを表し、「棒」（rupa）、「月」（indu sasīna）、「地球」（bhū）、「始まり」（adi）、あるいは「ブラフマ」など、何か特異なものを一般に指すもので表される。二は「双子」（yama）、「手」（kara）、「目」（nayana）などと表され、四は「海」（samudra）、五は「感覚」（viṣaya）または「矢」（カマデヴァの五矢）、六は「季節」または「味」、七は「山」（aga）などと表される。これらの名前は科学著作が詠まれた詩に適応し、また数字とは異なり容易な変更を許さず、そのため任何的改变都会扰乱韵律。

このシステムの例として、「['S]aka Sa[m.]vat, 867」（西暦945年または946年）が「\textit{giri-ra[s.]a-vasu}」によって与えられる。これは「山」（七）、「味」（六）、そして「 Vasu」の神々（八）を意味する。「日期を読むときには、これらは右から左へと読みます」という記述がある。

このシステムの発明時期については不明である。最初の痕跡はカティヤーヤナとラティヤーヤナの「['S]rautas[=u]tra」に見られる模様である［131］。ヴァルハーマヒラ（没年587年）には既知であり［132］、彼はそのシステムを使用した「B[r.]hat-Sa[m.]hit[=a]」がある［133］。また、['A]リヤバハタ（西暦500年代頃）がこのシステムを知っていたという主張もある［134］が、証拠は存在しない［135］。このシステムの最初の碑文的な例は604年と624年のベヤン（カンボジア）の碑文に見られる［136］。

また、南インドで発展した興味深いアルファベットによる数字符号についても触れなければならない。このシステムでは、次の表に示すように文字が数字を表す：

1        2       3       4       5        6       7      8     9      0
  k        kh      g       gh      [.n]     c       ch     j     jh     n
  [t.]     [t.]h   [d.]    [d.]h   [n.]     t       th     d     th     n
  p        ph      b       bh      m
  y        r       l       v       ['s]     [s.]    s      h      l

この計画では、一つの数字を複数の文字のいずれかで表すことができ、記憶のために単語に組み込むことが容易となる。例えば、

2     3      1     5    6       5       1
  $kha$ $gont$ $yan$ $me$ \textit{[s.]a$ $m[=a]} $pa$

という単語が形成される。

位取り記法を伴う後期のインド数字において、右から左に読み上げると156万5千1百3十二である［138］。これは、紀元後1184年に知られる最も古い表記であり、吠陀経義疏においてカリヤ Yug 以来経過した日数を表すものとされる［139］。バーンルは、このシステムが占星術の計算規則や天文台用の表を作成する際の記憶助具として現在も使用されていることを述べている。

また、類似の第二のシステムは、現在シーラン、タイ、ビルマにおいて手稿のページ番号に使用されている［140］。このシステムでは、34の母音を伴う34個の元音（ka...la）が1から34までの数字を表し、等価な35から68までの数字を[k=aa]...[l=aa]で表す［同］。さらに、i（ki...li）は69から102までの数字を表す［同］。

如既述したように、これまで説明したインドの体系は、ギリシャ人やヘブライ人の他の多くの体系と比較して向上するものではなく、ゼロが存在しないため、十、百、千などの高位単位を一から九までの単位で使用する同一の記号で指定することは実用的ではありませんでした。言い換えれば、位置価値があるためにはさらなる改善が必要であり、その可能性は存在しませんでした。そこでナナガハ（41）の符号は、「千七百二十四」という言葉を「T 7, tw, 4」のような現代の記号で書き表す必要がありました。つまり、千位の7、十位の2（「二十」が本来「二つの十」を意味し、$ty$は十を示す）、そして単位の4が言葉通りに与えられ、単位（十、百など）の順序は位置によって与えられます。この体系を完成させるにはゼロのみが必要でしたが、ナナガハの彫刻からおそらく八世紀後になって初めてこの重要な符号が現れました。そして、その符号が広く知られるようになったのはさらに遅い時期でした。発明に至った人物は誰か、彼が住んでいた場所はどこか、さらにはそれが何世紀の出来事だったかも、おそらく永遠の謎でしょう（141）。古代の算盤の一つの形態が、あるインドの天文学者や数学者に空欄を表すシンボルを使用することの有用性を示唆した可能性があります。インド各地で高位の名前は異なる形式を取り、その順序も入れ替わることが確認されています（142）。しかし、単位の意味が読み方によって与えられたため、これらの変動は誤りには至らなかったのです。実際、この変動自体が空欄や欠如した単位を表す言葉の導入を必要とし、最終的にはその言葉に相当するゼロの符号が導入されることとなったのです。

インド人が書き読んだ形で、大きな数8,443,682,155 を考慮する。これと対比して、ギリシャ人やアラブ人はどのように読み解いたかを示す。

\textit{現代アメリカの読み方}, 8兆4千4百3億6千8百2万1千5百5.

$インド語$, 8パドマ、4ヴァヤルバダ、4コティ、3プラユタ、6ラクサ、8アユタ、2サハスラ、1サタ、5ダサナ、5。

\textit{アラビア語と早期ドイツ語}, 八千千千と四百千千と四十三千千、そして六百千と八十二千と百五十五（または五と五十）。

$ギリシャ語$, 84万万と4368万と2千155。

ヴォープケは指摘したように、このような数値の読み方は、インド人が採用した表記法が場所概念を強調しようとしていたことを示している。他の言語では、サンスクリット語以外は、数の命名において十進法システムを一貫して適用していない。ギリシャ語における万、アラブ語および現代の数法における千の導入は、実質的に十進法体系から離れることを意味する。したがって、百未満の数では、英語の十一と十二は他の-teenの数に調和していない一方で、十から二十までのすべての数の命名は二十以上の数の命名とは類推できない。我々の書かれたシステムに則れば、十一つ、十二つ、十三つ、といった表現が必要であるが、実際には二十一つ、二十二つと同様に二十一、二十二と表記される。サンスクリット語は一貫しており、単位数が十と百の前に置かれている。また、どの古代民族もインド人がそのようなまでに数法を進めなかった。



当 $aṇkapalli$ が十進位置記法を完成させたとき、第十の符号は「空-bindu」と呼ばれ、一般的には「空」（そら）と略された。Brockhausは、ヒンドゥー教徒がゼロという符号の発明に適していたことについて、その哲学や宗教によって見事に適合していると言ったことは適切である。何も無を大きな成績の中心としたこの作業は、ヒンドゥーの天才と完全に調和している。

一般には、この「空」が符号として使用されたのは500年頃からではないかと考えられているが、一部の著者はそれよりも早く使用されていたとする説もある。アーヤバッハ（Aryabhata）は一般的な根の計算法を提供しており、彼は十進記法を知っていた可能性があるものの、我々の数字符号に由来する文字を使用していなかった。また、彼はしばしば「空」について言及している。もし彼が符号を指していた場合、「空」という符号は500年頃まで遡ることができるが、もちろん彼が単なる何も無の概念を指していた可能性もある。

やや後になっても6世紀にVarāha-Mihira（ヴァラハミヒラ）が『Bṛhat Saṃhitā』（ブラハサミタ）という著作を著し、しばしば「空」を使用して数字符号について述べており、彼が特定の符号を指していたとする説がある。これにより、アーヤバッハも同様に行動していた可能性が高まる。

また、興味深い点として、問題に関連するものとしてVarāha-Mihiraは上記のように位置値を持つ単語システムを使用していたことも述べておくべきである。

次の世紀の初頭（620年頃）にバンナ（B\[a\]nna）が『サヴァダタ』について称賛し、その中で星を零点に比喩していることを述べた。これらの星は現代のアラビア数字システムにおける零点と同様である。一方、この時期にインド人がゼロの象形文字を知っていたという有力な証拠はないとする反対意見としては、700年頃にアラブ人はシン province を占領し、そこで数値を表す方法を学ぶ機会があったが、776年に完全なシステムを受け取った際にはそれを新しく見出したと捉えたという事実がある。このような証拠は結論的ではないものの、8世紀初頭にインドで完全なシステムが広く普及していなかった可能性を示唆している。

しかし、1700年にドイツを訪れた旅行者が小数法を耳にした経験がないことからも、数学の専門家たちはアーラビハラ（\[A\]ryabha[t.]a）の時代にはゼロを知っていたかもしれないが、商人や一般の人々はその新しさの意味を理解するのに長い時間がかかった可能性がある。全体として、証拠はインド西海岸が完全なシステムが最初に見られた地域であることを示唆している。上記のように、カンボジアでもインド同様、位置価を持つ数詞の痕跡は非常に早期から見られるが、絶対的に十進法の象形文字が必要ないという点では異なる。

最初の銘文による九つの符号とゼロを含む位取り記法の使用について、確固たる証拠は存在しない。コールブックは1807年に、古代の銅版地図簿に含まれる多くの偽造の可能性について警告した。この点については、フィートはインドの铭文学者の研究を詳しく論じており、これらの偽造が主に十一世紀末頃に行われたと主張している。コールブックはカイよりもより理性的な見解を示し、カイはこれらの偽造がインド仮説全体の正当性を損なう可能性があると主張する一方で、「しかし、そのような疑惑がある場合でも、銘文が作成された時期に近い歴史的事実を利用することは完全には代替されない。偽造者がその作成を信じさせる必要があるため、彼は歴史に従うよう制約され、既存の概念に適合する形で行動しなければならない。そして、その時代に支配された伝統は、関心のある時期から離れているほど正確である可能性が高い」と述べている。ブーヘラーは、シディサムヴァト346年（595年）のグージララ銅版铭文を最も古い位取り記法の使用例として挙げており、「この铭文中で用いられた符号は当時のアルファベット数値と場所に一致している」と述べている。ヴィンセント・ア・スミスは815年（サムヴァト872年）の石刻铭文を引用しており、キエルホーンはパターリ柱铭文のパラバルが917年（西カルナータカ語：861年）に dated しているものも提供している。スミスが引用した铭文はバンドアルカーによって提供され、その他の铭文はサカサムヴァト715年（798年）のもので、柱石に彫刻されている。キエルホーンはまた、カナウジのマヘンドラパラの銅版铭文2つを提供しており、ヴァルハビサムヴァト574年（893年）とヴィクラマサムヴァト956年（899年）である。750年以降にまで遡る铭文が存在すること自体は、このシステムが少なくとも1世紀以上古いことを示唆する。さらに発展する過程で示すように、数字符号がヨーロッパに導入された後2世紀以上が経過した後、硬貨や铭文に現れた。チバートは特定の使用例を引用することを必要としないとする一方で、590年以降から痕跡が見られるという点を述べている。「現在すべての文明国で使用されているこのシステムがインド発祥であることは疑いようがない；他のどの民族もその発見に主張する資格はなく、特に西アジア諸国のこのシステムの起源に関する参照は一様にインドを指している」と述べている。

ブーヘラー、キールホルム、スミス（V.A.）、バンダルカール、そしてチバートの証言と意見は最も慎重に考慮されるべきである。古代インドの銘文学における権威として、彼ら以外の者は存在しない。彼らの研究は世界中のインド学者によって受け入れられ、位取り記法の発展に関する彼らの一致した判断は、その起源が六世紀前後のインドであったという点において、反証的な有力な証拠が出るまで否認されない。

多くの初期の著者はインド数字の多様性について言及している。アルビルーニーはおそらく最初であり、ジョハネス・ヒスパレニスも注目に値するものである。彼は七と四の変形を示している。第49頁に位取り記法を使用した数字符号の一覧表を挿入する。この一覧表の主要な資料源はブーヘラーであり、彼の著作には見当たらないが、非常に興味深いいくつかの例が提示されている。

表中の['サラダ]形式では、円が1を、点が0を表す。これらは、カシミールアスヴァーヴェーダの手稿と本文から抽出されたものであり、その手稿は確実に四百年以上前のものである。同様の形態はテューリンゲン大学所蔵の手稿にも見られる。

比較のために現代のサンスクリット語およびアラビア数字も加えておく。

サンスクリット, [図表]
 アラビア, [図表]

49

位取り記法を使用した数字符号

五十一

*       *       *       *       *

五十一

（注：原著中、各章の始末には図版番号が記載されていましたが、本文として翻訳する際は省略しました。また、原著中の数値「」もそのまま表記しています。）

\clearpage

\chapter{第4章：記号「ゼロ」の起源と誕生}
\section*{THE SYMBOL ZERO}

第四章

ゼロの符号の起源と歴史

原始的なブラーミーおよびカ热搜［＝t.］ヒ文字がインドで出現した数世紀後、初めてゼロがその地に現れましたが、そのような文字はキリスト教以前の数世紀にバビロニア人によって使用されていました。この符号は［バビロニアのゼロ符号］または［バビロニアのゼロ符号］であり、計算には使われていなかったようです。また、何らかの位の単位が欠けている場合でも必ずしも出現せず、例えば180は［バビロニア文字180］と記され、三六十分と単位なしを意味します。一方で181は［バビロニア文字181］と書き、三六十分の一単位です。このバビロニアの符号の主な用途は分数にあり、60分の1、3600分の1などであり、ギリシャ語における［ギリシャ文字o］（意味は「何もない」）の使用と似ています。

注185：バビロニア人によるゼロ符号の使用。
注186：単位が欠けている場合でも必ずしもゼロ符号が出現しないこと。

インドにおいて最初の確実なゼロの出現は、サムワット933年（西暦876年）にガウリヨールで見られる銘文にある。そこで50の花輪が記述されている箇所では［ガウリヨール数字符号50］と表記され、270も同様に［ガウリヨール数字符号270］と書かれている。バカルマス手稿はおそらくこの以前に存在し、点やドットをゼロの符号として使用していたと考えられる。ベイリーはオカマンデルで発見されたバージャー・ラシャトゥーラ王による738年（西暦）の恩給文書においてゼロが用いられていることを報告している。また、不鮮明なグプタ時代の貨幣である707年（西暦897年）のものにもゼロが見られるが、ベイリーの研究の信頼性は疑問視されている。如実に述べると、数字符号が銘文や貨幣上に現れるのは、そのシステムの起源と書面による説明よりも遙かに後の時期であることが推測される。前述の期間からインド全域へ急速に広まり、南部地域では現在も旧来のシステムを維持しているタミル語圏やマレーシャ語圏の人々がいる。

底本：David Eugene Smith \& Louis Charles Karpinski『The Hindu-Arabic Numerals』
該当章：第4章：記号「ゼロ」の起源と誕生

【翻訳】
早期の銘文に見られる以外にも、数論に関する初期著作において零の概念に対する特殊な扱いは、この符号がインド由来であることを示すもう一つの指標となっている。ウジャイエンに住み、インド天文学の中心地とされるブラフマグプタは、7世紀初頭に算術著作において零の性質について明確な扱いを加えている。彼は符号について論じていないが、その扱いから、何かしらの方法で零はギリシャ語や他の古代数論において見られない特別な意味を持つことになったと推察できる。ブラフマグプタ以外にも、バッカラはより科学的な扱いを加えているが、一度だけ零による除算という未許可の自由さを認める。古代インド数理学の最新に発見された著作であるマハーヴィラカーヤマの『Ganita-Sara-Sangraha』（830年頃）は、位値記数法を使用していないにもかかわらず、零との計算に関する類似した議論を含んでいる。

ゼロの形を示唆したものは、もちろん推測に過ぎません。ヒンズー人は手稿における欠落部分を埋めるために点を使用していましたが、これは文の途中で休止を表すときのように我々が使用するように、より自然な象徴でした；そしてこの点はヒンズー人が最初に使用し、現在でも多くのアラブ人が使用している象徴です。また、我々のXのような十字もこの目的のために使用され、時折ゼロの象徴として見つかります。上記で述べたバカルマンス手稿においては、「śūnya」という言葉に点を用いて未知数と零を表しています。また、比例式における未知数を表すためにゼロの使用例として、1467年と1468年にエrfurt大学で行われたゴットフリート・ヴォラックの講義に関するラテン語手稿においても類似の使用が見られます。この用法は18世紀以降でも指摘されています。

小円は十を表すために用いられた鋸状の円からおそらく影響を受けたものである可能性がある。また、ポタメイウスが『大成』において性単位制における偶然の空白を示すために使用したオミクロンも、インドの著者に影響を与えたと考えられている。この記号はヨーロッパとアジアで広く用いられ、アラビアの天文学者アル・バット＝イ（没年929）もアルファベット的な数字符号に関連して類似の記号を使用した。アル・バット＝イが性単位制における分（または秒）の欠如を示すためにアラビア語の否定詞「ラ」を使用したことは、ナルリーノによって指摘されている。また、「サラダ」文字においてカシミールのアスワヴァーディにおいても小円は一を表すために用いられており、その書記は少なくとも400年以上前のものであることが特筆される。

ブラフカラ（1150年頃）は数の上に小円を使用して減算を示し、タートル語では余分な単語を丸で囲んで除去した。私たちが「四在窟」のスコアマークと呼ぶものも同様の源から派生しているかどうか興味深い。オーキレット（1130年頃）は師であるバッサの書簡においてゼロを表すためにギリシャ文字の「τ」を使用し、これは「テカ」という言葉の略称であり、我々が後に見るようにゼロの一つの名称であったにもかかわらず、「チフパラ」からも来ている可能性がある。オーキレットはまれに「円柱」を使用し、両方の形を「シフラ」の名前で呼んだ。同じ記号を使用したのはフランス・シグソート（1150年頃）である。ハイベルグは15世紀の東ローマギリシャ人において他の特殊な形状が使用されていたことを指摘している。本文から明らかに、これらの著者は新しいシステムの意味を理解していなかったことが推測できる。

インドで最初に使用されたドットは、後に小さな円が代わり現在も使われています。しかしアラブ人はその円を採用しなかったため、アルファベットシステムにおいて数五を表す文字と類似していたからです。873年（西暦）の手稿で最も古く知られるアラブのゼロはドットであり、ジャムシの要約算術書「982年（ヒジュラ紀元）（1575年）」においても両方の記号が同じ意味を持つ場合があります。これらの作品では数は[シンボル]で表されます。5を表す形は、いくつかの作品では[シンボル]や[シンボル]となり、エジプトでは[シンボル]が見られ、印刷体では[シンボル]が使用されることがあります。現在ではアラブ人はヨーロッパの影響を受けたときのみ通常のシステムを採用する際にゼロを使用します。

中国では最初に明確な零の痕跡は1247年（西暦）のチンスンの作品に見られます。その形はヒンドゥー語の円形であり、おそらく旅行者によって中国に持ち込まれたものです。

この極めて重要な記号の名前についても注意が必要で、今日でもその呼び方についてまだ不確定な点があります。私たちは現在「ゼロ」「ナウト」「シファール」と呼んでいますが、電話オペレーターはしばしば「O」を呼び、無学または粗忽者は「アウト」と呼ぶこともあります。このような不確実性を考えると、過去にどのように呼ばれていたか調査する価値があります。

すでに述べた通り、インド人はこれを「スンヤ」、「虚空」または「空」と呼んでいた。これがアラビア語に移行すると、「アスィフル」または「イスフル」となった。レオナルド・ディ・パシァ（1202年）がヒンドゥー数字符号について記述した際には、この文字を「ゼフィルウム」と呼んでいた。ギリシャ語とアラビア語の両方の影響を受けたマキシマル・プランデス（1330年）は、「チプフラー」と呼んだ。イタリア語で書かれた算術に関する著作において、フィオレンツァのヤコブ（1307年）はこれを「ゼウエルロ」と呼び、アレッサンドリアのジョヴァンニ・ディ・ダンティ（1370年）の算術では、「チェウエルロ」と表記した。また別の形として「ゼピロ」があり、これは「ゼフィルウム」から「ゼロ」へと進化の一歩を示していた。

当然、英語の「cipher」やフランス語の「chiffre」はアラビア語の「a[s.]-[s.]ifr」から派生したものであるが、いくつかの言語では一般の数字符号を指す意味にまで拡大されている。この用語の多義性は、「ciphering」という単語にも現れている。「ciphering」は計算や計算演算を意味する。ヨハネス・フスウィートは両方の意味でこの言葉を使用し、十番目の文字について「theca, circulus, cifra, figura nihili」と四つの名前を与えている。

フスウィートのこの主張は、その時代の多くの著者がサクロボスコの『アルゴリズム』（1250年頃）に従っている可能性が高い。彼はまたハルフォックスやホリウッドのジョンとして知られている。サクロボスコの『アルゴリズム・ vulgaris』に対するペトゥス・デ・ダーシアの解説（1291年頃）も広く使用された。

その時代の大学でこのイギリス人の算術に関する著作が広く使用されていたことは、13世紀から17世紀までの多くの手稿の存在から明らかである。ミュンヘンには20点、ウィーンには12点、エルファートには13点、イギリスではハリウェルによる記載がある数点、コックスの『オックスフォード大学図書館目録』にリストされている10点、プラムプトン収蔵品に1点、コロンビア大学図書館に1点、もちろんそれ以外にも多くの手稿が存在する。

从「a[s.]-[s.]ifr」发展而来，最终演变為「zephyr」「cipher」，最后固定為「zero」。この最終形態が印刷された最初の作品は1491年カランディの算術に見られるようである［227］。手稿では、少なくとも14世紀中頃から現れている［228］。また、1485年にジェハン・セレンティン著『Le Kadran des marchans』にもこの語が見られる［229］。

この言葉は比較的早くスペイン［230］とフランス［231］で広く知られるようになった。中世の作家たちは、しばしばこれを「sipos」と呼んだ［232］。また、「wheel」「circulus」（ドイツ語では「das Ringlein」［235］）や「circular note」「theca」などと呼ぶこともあった［236］［237］。この名称は、ギリシャ語のthetaに似ていることから「null」とも呼ばれていたが［240］、ペトロス・デ・ダーシアによれば、これは盗人や強盗に額や頬に円形の印をつけるために使われる鉄の名前から派生したと説明されている［238］。また、「omicron」（ギリシャ語の「o」）とも呼ばれ、しばしば「o」または「ph」で区別するために書かれた［239］。さらに、「nihil」や「nulla」（ラテン語の本では）［241］［242］や「rien」（フランス語では）［243］、そして最も一般的には「cipher」と呼ばれていた［244］。

ウォリスは［245］、「zero」に関する様々な形態について最初の延べした議論を提供しており、その他の重要な変種として「ziphra」「zifera」「siphra」「ciphra」「tsiphra」「tziphra」とギリシャ語の「tziphra」も挙げている［246］。

\clearpage

\chapter{第5章：ボエティウスによるヨーロッパへの数字導入をめぐる問題}
\section*{THE QUESTION OF THE INTRODUCTION OF THE NUMERALS INTO EUROPE BY BOETHIUS}

第五章

ボエティウスによる数字符号のヨーロッパへの導入に関する問題

私たちが数字符号の起源について確定的な結論を下すことができないかのように、その導入時期と場所についても確信的な判断は下せません。この導入には二つの一般的な説があります。一つ目は、モーセによって8世紀または9世紀にスペインへ持ち込まれた後、キリスト教ヨーロッパへ伝えられたという説で、これについては後述します。もう一つはWoepckeが唱えた説であり、モーセによってスペインへ持ち込まれなかったというものです。アルガズールの到来時にはすでにスペインに存在していたとし、その数字符号はネオピタゴラセス経由で西方に伝播したとするのです。

この二つの説を支持する事実が二つあります。（1）これらの数字符号の形態は特徴的であり、13世紀初頭（1202年以前）にレオナルド・ダ・パルマからノルthern Africaへ持ち込まれたものとは著しく異なる。（2）これらは伝統的にボエティウス（500年前後）に帰せられる数字符号であり、彼がこれらのネオピタゴラセスやその源流から直接受け取ったとしても不自然ではありません。

さらにWoepckeは、711年アルガズールがスペインに入国した際には、税金計算のための数値システムを征服地のものに採用するという習慣があったことを指摘します。そのため、スペインからイタリアへ持ち込まれた数字符号は東方アルガズール帝国のものとは異なる変化を経験せず、バグダッドを通じて伝播したものと異なる形態を取ったはずです。

この説によれば、ゼロを含まないヒンディー系統が450年頃アレクサンドリアに到達し、神秘的な特に東方への愛好心からネオピタゴラセスはその使用を奇妙なものとして受け入れたとされます。その後地中海を経由してアテネやローマのボエティウスへ、そしてスペインの学校へと伝播し、アルガズール自身が改善されたシステムを知る以前にアフリカやスペインで発見されたというのです。

ボエティウスによるヨーロッパへの数字符号導入に関する最近の理論は、Bubnovによって提唱され、その内容が慎重な意図から成るため、特に注目されるべきである。Bubnovは、ヨーロッパで最初に見出された数字符号形態が算板上使用された古代の記号から派生したものであり、しかしゼロはヒンディー語源であると主張している。しかしながら、既に提示された証拠からはこの理論が支持されない。

ウープツェックの理論によれば、二つの問題が提起される：（1）これらのスペイン数字符号の性質は何であり、どのようにイタリアへ知られたか？（2）ボエティウスはこれらを知っていたか？

スペイン語での数字符号は「\textit{ huruf al-ghubār}」と呼ばれ、「ghubār」または塵の数字符号と区別される「\textit{ huruf al-jumal}」やアルファベット的な数字符号とは異なる。おそらく後者はシリア人やユダヤ人の影響によりアラブ人も使用していた可能性がある。語彙「ghubār」の意味は、これらの数字符号が塵の算板上に記されたものであることを示唆していると考えられる。これは数字を表すためのカウンター法とは異なる方法である。また、アルビルーニによれば、ヒンディー人たちはしばしば砂で計算を行っていたという点も注目に値する。この用語は950年頃からキアールワーンの匿名作者による詩歌の中で見られ、「ghubār」計算に関する彼の著作について述べている。その後、アブ・ベクル・ムハンマド・イブン・アッバド、通称アルハサールという名前のアラブ人作家は、その著作の二章目を「塵の数字」についてとし、この理論が支持されている。

ゴバール数字符号自体は、シルベストレ・デ・サシーによって初めて現代の学術界に知られ、古代サン＝ジェルマン＝ド＝プレ修道院の蔵書からアラビア語手稿で発見された。このシステムには九つの文字があり、ゼロは存在しない。一桁の上に点が打たれると十を表し、二つの点は百を表すなど、[5・点]は五十を、[5・三点]は五千を意味する。これらの点が数字符号の上に散らばったことからゴバールという言葉が生じた可能性があると指摘されているが、これは必ずしも確実な説ではない。この点のシステムは、後のペルシャで見られ、現代アラビア語に近い数字符号とともに見られるが、その使用自体は重要な意味を持つ。なぜなら、完備されたインド数字符号が近くにあるとき、西洋のシステムがペルシャに戻る可能性は低いからである。

最初の印象では、ゴバール系統におけるこの特徴がアラビア語由来であり、現在の彼らのゼロである点がこれに由来しているという理由があるように思われる。セミティック民族がこのような計画を採用したのは自然なことである。なぜなら、その文字の発音符号がそれを示唆し、またギリシャ語数字符号におけるグリークアクセントの影響も考慮に入れれば然りである。しかし、バカシール手稿において点がゼロとして見られることを考慮すると、十世紀の『キタブ・アル＝フィヒルスト』において下付き形式で使用され、十六世紀までその使用が続いていることから、この形態もインド語由来である可能性がある。

アラブ人は、既に述べたように、ヒンズー零をすぐに採用しなかったのは、その形が5に似ていたためである。彼らは上付き点を使用し、これで十分な目的を果たしていた。実際には、この上付き点はアラビアとヘブライの学術者たちによって10世紀の北アフリカでインド数字符号として知られ、アブ・サル・イbnタミンによる『セフェル・ユシヤ』の解説においても「ゴバールまたは塵計算」としてその名が使われている。これらの事実から推測すると、773年に再びアラブに伝わる数字符号以前にも、アラブ人はゴバールの形を知っていた可能性が高い。また、「ゴバール数字符号」の用語は、上付き点の特殊な使用法に関連するものではなく、その名称で使われていた。この文脈において、アルジェリアでは14世紀後半までの手稿においても異なる2種類の数字符号が使用されており、現在のアラブ文字ではなくヨーロッパの形式を使用しているモロッコ人は今日でもその傾向を継承している。

インドでは十、百、千等を示すために下付き点を使用していたという事実は、987年頃の『キタブ・アルフィヒルスト』に記された一文によって確認される。この文中で作者はインドの人々の書記法について議論している。この引用は数字符号の初期歴史にとって重要な参考文献でありながら、以前の著者はこれを指摘していなかった。ここで示された数字符号は通常「インド」と呼ばれるものであり、ゴバールとは対照的である。この文書では、点が文字の下に配置されており、先述した上付き点とは異なる意味を持つものであった。

ボエティウスによる数字符号のヨーロッパへの導入について、その形式は我々のものと比べて、アラビアの数字符号よりも類似性が高い。これらは多少異なるが、大まかには以下のようであった：

1[図]
  2[図]
  3[図]
  4[図]
  5[図]
  6[図]

エジプトのデモティックとヘリオティックの序数形式が可能な影響を与えた可能性はしばしば指摘されており、この点を比較するためにこれらの形式を紹介することが適当である。これらの形式はヒンドゥー語の形式とも関連があり、最初の三つがこれらすべてのシステムと関係があることは明白である。さらに類似性のあるものはデモティック4と9のみであり、したがって、ヒンドゥー語の形式がこの源から来たという主張には基礎がない。エジプトの初期の四つの基数数字符号は現代のアラビア語に近い。

[図：デモティックとヘリオティックの序数]

ヨーロッパへの極めて早期の数字符号導入説は、いくつかの点で失敗している。まず第一に、初期の西欧形式は知られていない；第二に、初期の東洋形式が [.g]ob[=a]r と似ており、69ページの第三行から Shiraz の手稿（約970年頃）の形式を見ると、一部の西アラビア形式も使用されていることが確認できる。最も重要な事実は、Sacyによって示された [.g]ob[=a]r 数字符号が、八を除き、単一のアラビア文字または文字の組み合わせである点である。Woepcke理論と [.g]ob[=a]r 数字符号の意味について述べた以上に、ボエティウスはこれらの [.g]ob[=a]r 形式を知っていたか、それらと類似した形式を知っていたかという問題を考える必要がある。

ボエティウスとは誰か。中世には「Divus ボエティウス」と呼ばれた アニキウス・マヌリウス・セベリンゥス・ボエティウスは、475年頃ローマで生まれた。彼は安シイ家族のメンバーであり、その家族は彼が生まれる前から数十年間クリスチャンだった。幼少期に父を失い、伝説によれば十歳頃にアテネへと連れて行き、十八年間在那里学习。ローマ senator シマキウスの娘 ルスティチアナと結婚し、両家の力強い結合により彼は最高級の社会で活動することが可能となった。彼は厳格に正義を主張し、宮廷内の不正に対して反対したため、王座近くの無恥な要素から憎しみの的となり、元副将アルビヌスが不当に叛逆罪で告発されたことを勇敢に擁護した結果、ピアヴェで拘禁され524年に処刑された。彼の死後数世代には至らず、歴史的な批判が最も低かった時期に教会は彼の処刑を殉道として利益を得たと見做すようになった。その結果、彼は聖人として扱われ、骨も安置され、自然にその書物は千年間教会学校における古典文献の一つとなり続けた。しかし、ボエティウスの算術が抽象的で意味不明であり、必要以上に複雑であることを考えると、中世の学生がこの作品から精神的な栄養を摂取しようと試みる姿は悲劇的である。

彼は同時代人および直後の後継者たちから、大師として評価された。カシオドロウス（490年頃 - 585年頃）は彼に対して次のように述べている。「あなたによって翻訳されたピタゴラスの音楽とポタメイオスの天文学はイタリアの人々に読み解かれ、ニコマコスの算術とユークリッドの幾何学は西欧の人々に知られている」と述べている。中世シュolasticismの創始者として、三科と四科を区別し、彼の時代唯一の古典を著したギボンは彼を「カトゥスやテルレウスが祖国人として認められた最後のローマ人」と呼んだ。

ボエティウスに関連する第二の問題は、彼がヒンドゥー数字符号を知っていた可能性があるかどうかである。東洋と西洋との関係性が示されることにより、この質問に対する肯定的な答えしか存在しない。数字符号はゼロを含まずに数世紀間存在し、インドではよく知られており、東西間には継続的な思想交流があった。また、戦士、大使、学者、そして不休の商人が地中海地域とインドの貿易・文化中心地間で陸路や海路を通じて行き来していた。ボエティウスは旅行者や商人からヒンドゥー数字符号の一つまたは複数の形式を学ぶことができたと考えられる。

この主張を正当化するためには、遠東とヨーロッパとの関係についてより詳しく述べる必要がある。実際、ボエティウスの時代以前の一世紀においては、東アジアと中央ヨーロッパ間の大規模な学術交流に関する記録は存在しない。しかし、学者たちが知識の唯一の伝播者であるという誤りを否定する理由がある。事実上、ヒンディー語数字符号がアラブに知られるべき自然な流れであり、ゼロが場所値を可能にするまで遙か以前から、これらの商路を通じて遠西まで広く知られていたと考えるのが適切である。また、その文字、計算方法、時々の改良、ゼロの登場、そして商業問題解決に関する習慣も、東洋から西洋へと世代を超えて伝わっていた。

学問的な知識の拡散は、商人や浪人学者に依存していたことを常に念頭に置いておくべきである。実際、アビシーナ（980-1037）自身の短い自伝によれば、彼の家族がボカラに住んでいたとき、父親は彼を果物屋の家に送り、ヒンディー語計算法を学ばせた。この果物屋（油商かもしれない）はその分野で専門家だったという。ピサのレオナルドも同様の訓練を受けた。

商業知識が商路を通じて広がる全体的な問題は、ゴバール数字符号、ボエティウスの問題、ゲルベルト、ピサのレオナルドなどに関連しており、これらの名前や出来事について考慮するための一時的な逸話が必要となる。

六世紀前後において、ナナガット洞窟にヒンドゥー数字符号が彫刻される以前から、アラブとインドには貿易関係があった。実際、アーリアン族がインドへ移住するずっと前に、大きなトランシルビア族は地中海からインダス川まで文明を広げていた［291］。その後の時代において、アラブ人は西側のエジプトとシリアとの間、そして遠東との間に中継役を果たしていた［292］。六世紀前半には、地理学の父ヘカテウスは地中海沿岸国だけでなく、インダス川まで到達する国の知識も持っていた［294］。また、旧約時代には三年に一度のインドへの航路が定められていた［296］。ヨセフの物語は、インドからアラブを経てナイル河畔までの caravan 貿易を証言している。

スキュラクとほぼ同時期に、ダリウスの海軍将校カリアンダ海岸出身のスクリラックスがインディアン西北部へ旅し、その航海記録を著した［295］。彼はインダス川沿岸諸国をペルシャの支配下に引き込み、当該地域で存在していた数制度も当然ながら彼の能力に照らし合わせて知られていた。

スキュラクから百年後には、ヒロドトゥスはインドについて相当な知識を持っていた。その中でインド産の綿や金を述べ［296］、セソトリスがインドへ航海する船を編成したことを語り、東方への路線も挙げている［297］。これらの路線は一般に紅海経由であり、フェニキア人やサバエ人がそのルートを使用し、その後の時代にはギリシャ人とローマ人もこのルートを採用していた［298］。

第四世紀前半には、東西の関係は非常に密接なものとなりました。330年頃には、マッソリウム（マルセイユ）出身のピテアスが、ブリテン諸島の北端およびドイツ海沿岸まで探索を行い、またマケドニアは南部フランスと緊密な関係を保ちながら、アレクサンドロス大王の軍隊はアフガニスタンを通じてパンジャブまで遠征しました。プラinyによれば、アレクサンドロス大王はインドの道路測量に地図製作者を雇用したと伝えられており、メガステネスは前295年、セレウコスの大使としてパタリプトラ（現在のパトナ）に滞在し、その偉大な道路の一つについて記述しています。

また、ヒンドゥー人はギリシャ人からまたはおそらく中国から貨幣造作術を学び、現在北部インドで見つかるギリシャ・ヒンドゥーの貨幣は歴史的情報を提供する重要な源泉となっています。ラマーヤナには、多くの商隊が旅し、海路で外国へ出航することについても言及されています。セイロンはマラッカとシアムと貿易を行い、ジャワはヒンドゥー商人によって植民され、このように数世紀にわたる数字の形成期には商業知識がインド諸島全体に広がっていました。

また、早期のギリシャ人侵攻の結果は、マッセナのディカエルコス（紀元前320年頃）によって地図にまとめられ、長く標準的なものとして使用された。さらにアレクサンドロスは東への影響を断ち切らなかった。インドを三つのsatrapieに分割し、二つにはギリシャの総督を置き、残り一つにはヒンドゥーの統治者を残した。またバクトリアナ、アリアナまたは古代ペルシャの一地区には総督を置いた。これらの地で西洋文明は長く存続した。この時期にギリシャやローマの音律や天文術用語が Sanskrit に取り入れられたことは疑いない。そして紀元後2世紀から5世紀まで、インドの貨幣にはヘレニズムの影響が見られる。ヒンドゥーの天文術用語は西洋思想との関係を示しており、ヴァラハミーニャ（6世紀）とアーラビヤータ contemporaneous に、彼の著作『ブラハサミッタ』は「大作」の直訳である『プラトメイの megale suntaxis』と名付けられた。このように思想の交換は様々な形で明らかになっている。古代ギリシャ人がインドを訪問することは全く不自然ではなかった。ストラボはそのルートを明記しており、すべての旅人はエフェソスから出発し、フリッジアとカパドキアを経て直線的な道を進むという。東からの商品は常に西洋に輸入され、ギリシャ人はマラバールから生姜を得た一方で、フェノキシengeanceはかつてマラッカから金を運んできた。

ギリシャも中国との早期の関係があったと考えられる。その理由は、ギリシャと中国の生活様式には著しい類似性があり、これは彼らの建築様式、家庭的な習慣、結婚儀礼、公の物語者、ヘロドトスがエジプトから導入されたと述べている提灯劇、通りの占い師、サイコロゲーム、指で数を当てる遊び、水時計、音楽体系、万卾の使用、暦法など、多くの面で見ることができるからである。今日、大鐘寺へ向かう途中の北京郊外を通過していると、古代地中海文明に近い半ギリシャ風建築が連想されるほど、現代の中国は古典的な文明と密接に関わっている。中国の歴史家によれば、紀元前200年頃には彼らの武力が遠西で成功し、紀元前180年には大使がギリシャ都市バクトリアに派遣され、そこで中国製品が市場で販売されていることを報告したという。また、特定のギリシャ語と中国語には著しい類似性があり、これは遙かな昔から思想的な交流があったことを示唆している。

ローマ人はまた東方と商品を交換していました。ホラセウスは言います：「活発な商人よ、遠くのインドへ急ぐべし、貧困から海を越え岩場を越え火を越えて飛ぶべし。」 オrientの製品、インドからの香料と宝石、ペルシャからのフランキンセンス、中国からの絹は地中海諸国より需要が多かったため、貿易収支は西に不利となり、結果としてローマの貨幣は東へと流れていきました。例えば1898年には、パクルイ[=i]で発見されたローマの貨幣群があり、その年代は紀元前114年からハドリアヌス時代までであり、アボタバード北十六マイルにあるハザラ地区の一部でした。また時々類似の発見がなされてきました。

アウグストゥスはインドからの使者を受け入れたことを述べており、これは以前に知られていなかったことでした。彼も中国から大使を受け入れた可能性もあります。シュートゥイウス（紀元後一世紀）はその歴史の中でこれらの関係について言及しています。また、彼の同時代者数名も同様に述べており、 Vergiliusはアウグストゥスがペルシャで戦ったことを語っています。プラーニウスの時代にはローマ帝国とアジアとの貿易額は年間一千五百万ドルに達し、当時の相対的な数値よりも現在の方が遥かに大きかった。コンスタンティヌスの時代にはヨーロッパは遠東と直接通信を保っていた。

これらの関係を鑑みれば、ギリシャ人やローマ人がインドの数制度について何か知っていたという証拠がいつか明らかになる可能性があると考えられます。幾つかの著者がそう主張しているように。

東アジアに戻ってみれば、インド自身における知識の拡散とその周囲に関する証拠は多く存在する。紀元前三世紀に仏教が思想の媒介として発展し始めた。すでにヒマラヤ地域を通過しており、東トルキスタンまで到達していた。おそらくその後、中国へも伝播した。数世紀後の六二年（西暦）に中国皇帝はインドへの大使を派遣し、六七年には仏教僧侶が中国に招致された。また、この時期にはインドのアーソク王が既に本作で言及したように、領土をアフガニスタンまで拡大しており、パنجاب地方の数字符号がその早期から北へ伝播していた可能性も十分にある。

さらに、ペルシャの影響は知識の伝播において忘れてはならない。五世紀にはジョンディ・サプールのペルシャ医学校ではヒンドゥーとギリシアの学説が共存しており、ファルデス＝イ告诉我達，在科斯ロ伊斯一世的辉煌统治时期，即波斯巴列维文学的黄金时代（八二年），印度象棋被引入了佩尔西亚。当时正值与希腊的战争使萨珊王朝声望大增之时。

また、ボエティウスの時代から六世紀にかけて、エジプトの僧侣コスマスは trader の初期にアビシニアやインド、セイロンへと旅し、『インド航海者』という名を受ける。彼の地図（五四七年）には大西洋からインドまでの地球に関する知識が示されている。このような人物であれば、滞在した人々が使用する数字符号を観察することはほぼ確実であり、その研究結果を記録せずとも口頭で伝播させる可能性が高い。

アラブにおいては、その活動がムハンマドから始まったという認識は誤りである。彼らは商取引を常に尊敬の念を持っており、コレイシュは多くの世代にわたって年間を通じてヤマン地方の香料や織物を地中海沿岸まで運ぶ caravan を送っていた。5世紀には、アラブ海でインドとさえ中国との商取引も行われており、ヒラは東洋製品の交易拠点であったため、世界中のどの数字符号システムも孤立することは稀であった。

戦闘的な活動が始まる以前からアレクサンドリアは世界的な市場の中心地となっていた。この中心地から caravan がアラブを経てハドラムートに至り、そこでインドからの船と遭遇した。また他の caravan は北へダマスカスに向かい、さらに別の caravan は地中海沿岸南部を進んだ。シナイ半島の地峡から南方ヨーロッパのすべての商業港湾までそして黒海に至るまで、船舶が航行していた。アレクサンドリアのバザールを訪れる商人達にはヒンドゥー人がおり、ベロニカスもまたビザンティウムに報告されている。

SUCHはボエティウスの時代にアレクサンドリアの商人や地中海の他の海港都市の商人がパंजाबやインドその他の地域、さらには中国やさらに遠くのペルシャ、シヤンやマレー半島の数字符号を知っていたことを示す証拠の簡潔なまとめである。ブラーマ符号は学術的な注意を集めなかったかもしれない。なぜなら、ゼロが存在しなかったためで、他の多くのシステムと比べて優れていたわけでも劣っていたわけでもないからである。もしボエティウスがこれらの符号に興味を持ったとしても、それは自然に不思議なことや神秘的なものに対する興味と同じようにであり、彼はそのような情報についての興味深い片鱗を自分の著作でしか偶々述べただけであろう。

したがって第二の質問に対する答えは、ボエティウスがヒンドゥー符号を知っていたかどうかである。この問いに対する回答は、絶対的な疑いなく、彼がそれらを知っていた可能性が高いと断言できる。彼のような好奇心旺盛な人物がこのような興味深い情報を拾わなかったとは考えにくいからである。

では第三の質問に移り、ボエティウスがこれらの符号を知っていたかどうかの具体的または有力な証拠は存在するかを検討する。この問いは新しいものではなく、かつて二百数十年前にウォリス（1693年）がヴォッシウス（1658年）によって指摘されたことについて疑問を呈した後も未だに完全には解決されていない。ボエティウスと関連付けられた数学に関する三つの著作がある：（1）算術、（2）音楽に関する論文、（3）幾何学である。

算術および音楽に関する論述の真偽は、一般に認められており、しかし幾何学においては、ゼロを含むヒンディー語数字が含まれている部分について疑惑がある。ボエティウスがこれらの数字を知っていたことや、この書籍に包含していたという主張は多数存在する一方で、これらの数字を含む部分は偽物であると感じる見解も同程度の支持を得ている。特に疑問視される論点に対する否定的な立場を支持する者の主張は次のように要約できる。

1. 文献改ざんは常に問題となってきた。ローマ時代においても、中世においても、そして今日においても、その頻度は我々が想像するよりも遥かに高い。私たちは、どの聖歌集の編集者が自らの裁量で古典文学を任意に変更することを正当化されるかのように、古典文学の改ざんがいかに容易であるかを見るだけで、中世の抄録家が文書の段落を追加または削除するのにどれほど容易であったかを理解できる。

2. ボエティウスがこれらの数字を知っていたとしても、算術においてそのことを明記していたはずだが、彼はそうしなかった。

3. 仮に彼がそれらの数字を知っていたとしても、他の著作でそれを言及していたとしても、同時代者や弟子たち、さらには後継者がそれに触れるべきであった。しかし、カペラ（475年頃） やボエティウスの著作を知る中世の多数の著者は、そのシステムについての言及を見つけることはできない。

問われた文書は、ある抄録者が加えたものに見られる特徴を有している。ボエティウスは幾何学の著作において角度について述べているところに至ると、突然数類の議論へと文章が転じる。その後、算盤に関する説明の章が続き、その中に呼ばれる「アピクス」や「キャラクタレ」という数字符号の形式が記述されている。これらの符号の形は手稿によって異なるが、一般的には第88頁に示すような形である。これらは通常、左から9を始め右へ単位まで減少する形で書かれ、多くの著者はこの順序が反西欧的なシメティック語源からの影響によるものだと述べている。この慣習は16世紀まで続いた。その後、作者は主題を完全に離れて、残りの議論ではローマ数字を使用する。ボエティウスの方法とは全く異なるこの手口は、真実性の仮説に基づいて説明が困難である。なぜこのような学術的な著者がその起源や使用法について言及せずにこれらの数字符号を提供したのか？彼はそれらに関連する歴史的興味を述べるか、あるいは議論の中で使用するか、絶対に現在のままの文書を残すことはなかったであろう。

ボエティウスがヨーロッパに数字符号を導入したかどうかについて、サーE・クリブベイリーはさらに証拠として挙げているのが、4がナガートタイプではなくカブール形式である点である。アラブ人は776年までその形式を受け取らなかったため、ボエティウスの時代にヨーロッパでこれらの文字が認識されていたとしても、この特定の形態は認識されなかっただろう。また、フリードレインによって示された主要手稿から得られる六つの算盤符号においても、ゼロのいずれかの形式が含まれており、その象徴はおそらくこの時期またはそれ以降にインドで発生したと考えられる。ヨーロッパにはそう速く到達する可能性は低いと思われる。

第四の問題について、ボエティウスが数字符号を知っていたかどうかである。現在の証拠によれば、（1）ゼロがない場合でもボエティウスがこれらの数字符号を知っていたことは容易に考えられるが、（2）彼がそれらを知っていた確かな証拠は存在しない。ボエティウスが接触した可能性があるように、他の好奇心旺盛な心もその時代または10世紀に確定的に現れる以前のどの時代でも接触することができた。この5つの世紀は最も深い中世 darkest of the Dark Ages であり、これらの数字符号が偶然アルクインの著作に入り込んだ場合を除き、その痕跡は当時の文献には現れにくい。実際、これらの数字符号の存在に関する具体的な証拠は9世紀または10世紀から初めて見られるようになった。おそらく各地の商人が知っていたが、不完全な状態では重要な影響力を持たず、大きな注目を集めることはなかった。

この簡略な証拠の調査から、以下の結論が推測できる：（1）東と西との間には、商取引や旅行という名目以外でも、絶えず交流があった；（2）商人たちは交易相手の原始的な数の体系について知る機会があり、また無理に考えれば彼らはその知識を永久的に記録することはなかった；（3）浪人学者たちは出会った人々に関する多くの奇妙な知識を持っていたが、作家としての傾向は希薄だった；（4）アラブ軍が北アフリカを征服する以前から、商取引においてヒンディー・アラビア数字が知られていることは理にかなっているし、おそらく一部の浪人学者にもその知識があったと考えられる；（5）10世紀後半にこれらの数字符号が知られていたこと、そして1202年にフィボナッチによって詳細な著作が執筆されたこと自体は驚くべきことではなく、それ以前に現れたより広範な手稿証拠が見つからないことに驚きを感じるべきである。中世の初めから、少なくとも多くのキリスト教徒商人たち、おそらくいくつかの学者にはゼロを含まないこれらの数字が多少は知られていたことは疑いようがない。その文書的な証拠の欠如は、すべての状況を考慮に入れると全く不思議なことではない。

\clearpage

\chapter{第6章：アラビア世界における数字の発展と受容}
\section*{THE DEVELOPMENT OF THE NUMERALS AMONG THE ARABS}

第六章

アラビア世界における数字の発展と受容

仮に数字がインドで起源をもつとするなら、アラブ人がそれらをどのようにして知ったかはいつ頃のことであったろうか。モーレハムマディanismの影響により学問がペルシャやアラビアへと広まり、それが主な理由であるという説は一般的である。部分的にはその通りかもしれないが、モーレハムマディが出現する以前からこれらの地域では学問が既に尊敬され、商業もこの地域全体で盛んに行われていた。例えばペルシャにおいて、キスロヌシルワーン王の在位期はジュスティニアヌス法皇と並行して偉大な時代であり、社会経済状況の改善だけでなく、学問の振興も特徴としていた。キスロヌシルワーンは学問を奨励し、ギリシャから学者を招き入れ、西からの文化とともに東からの文化的影響を促進した。この時期にはアリストテレスとプラトンの著作が翻訳され、サンスクリット語からペルシャ語への『ヒトパデサ』やピルパイの寓話の一部も翻訳された。これらは、バグダッドにおける偉大な知的支配以前に三世紀以上前からペルシャで学問の振興が行われていたことを意味し、モーレハムマディ出現以前の影響のもとで進行していたことを示唆している。

第七七三年に、インドの天文学者が哈里发の御所を訪問し、天文学表を持参したことが、ヒンゥー数系統がアラブ世界に導入された最初の明確な痕跡とされる。これらの天文学表は、ハ里發の命によりアルファザーリによってアラビア語に翻訳された。アルカワルズミ（アブドルカフタール）およびハバシュ（アヒメド・イブン・アブダッラー、八七〇年頃没）は、アルファザーリの著作に基づいて有名な天文学表を作成した。ヒンゥー数が同様に導入された可能性が高いと断言できる。それらは、アルカワルズミの時代から数十年後には既に知られており、八二五年頃に『アルゴリティミス・デ・ヌムロ・インドオーム』の原本が著される前から、ヒンダーニュマークルに関する最初の著作ではないことを主張している。

上記の三人の著者は、第八世紀末から第九世紀末までの期間をカバーする。歴史家アルマサディおよびアルビリヌイは、前述した人物たちに近い時期に活動したが、アルカワルズミと同時期にヒンダーニュ算術に関するアラブの著者についても再確認すべきである。第一章で述べたアルコンディ、サンディ・イブン・アリ、アルスフィーである。

五百年以上にわたって、アラビアの著者たちはまた他の人々も、算術に関する著作を「インドの」という名称で呼んでいた。10世紀には、アンティオキアの`Abdallāh ibn al-Hasan、Abū 'l-Qāsim（987年没）およびマフムード ibn `Abdallah、Abū Nasr（大約982年）といった著者がいた。同じ時代やそれ以前に『フィヒルスト』（10世紀末）において言及されている者には、ハーラーンのSinān ibn al-Fatḥおよびアフマド ibn `Umar、アルカビシ（アラビア語表記：アルカバシ）が含まれる。11世紀にはアルビルーニ（973-1048）、アリ ibn アハメド、Abū 'l-Hasan、アルナサウィ（大約1030年）といった著者がいた。その次の時代にはイシャク ibn ユスフ・アルサダフィ（大約1174年）およびサムール ibn ヤヒヤ・ibn アッバス・アルマリブィ・アルアンダルシ（大約1174年）の著書が見られる。そして13世紀には、`Abdallātīf ibn ユスフ ibn マフムード、ムワファック アルディン Abū マフムード・アルバガダィ（大約1231年）およびイブン アルバンナの著書がある。

ギリシャの僧侣マキシマル・プランウデスは、14世紀前半にその著作『インド算術』においてアラビア語の用法を踏襲した。算術に関する多くの他のアラビア語の著者も存在したが、その多くは数字符号の起源について言及する必要を感じなかった。これらの符号の知識はおそらく当然とされており、その起源を明記することが必ずしも必要ではなかったのである。

九七〇年頃の文書には、ゴバールとアラビア語の通常の数字形が共存していることが示されており、これは特別な興味を引きます。この作品のタイトルは「数に関する興味深い問題集」で、アル＝ムハンマド・イブン＝アブダルジャリル・アル＝シージジ（九五一年～一〇二四年）が、ナジフ・イbnユムーン（アル＝カッサス）の著作から抄録したものです。ナジフのこの作品については、スーテルは言及していません。

第二に、純粋なアラビア語の影響を過大評価しない理由の一つは、アラブ自身がいかなる知的な強さも示さなかったことです。ムハンマドが人々の心に火をつけた後に行われたことは、異なる種類の強い民族が融合したときによく起こる現象と同様でした—多様な分野における大きな再生です。そのような種類の融合はテレスス時代のミレトスにおいて、ローマ初期の侵略者時代に、ギリシャ人がエジプト土壌に定着したアレクサンドリアにおいて見られ、現在では新世界ですべての民族が活力を混ぜ合わせるときにも見られます。したがって、アラブ文化とペルシア文化が融合した際には、新しい文明が生まれ、栄えた。アラブの影響は純粋さからではなく、より文化的ではあるが力強さに欠ける他の影響との混ざり合いから来ました。

第八世紀から第十三世紀までの間、イスラム教は世界に与えた影響は、ローマやアテネ、そしてイタロ・ヘレニズムの影響が古代文明に与えたものと同等であった。「彼らにはギリシャ人やヒンドゥー人のような発明心はなく、中国の天文学者のように観測に持久力を持たなかったとしても、少なくとも新しい勝利を収めた人々としての生命力があり、他者の成果を理解しようとする意欲と、数学や詩、哲学、言語学への同等の熱意があった。」

六二二年にムハンマドはメカから亡命し、その約百年後には十字架がキリスト教アジア、北アフリカ、そしてスペインの大半で新月に取って代わられた。アラブ帝国は学問の楕円形を形成し、その焦点はバグダッドとコルドovaであり、統治者はしばしば征服の結果として知的財産ではなく商業財産を求めることに自慢した。

これらの影響の下で、仮にムハンマド以前からもしくはその後において、ヒンドゥー数字符号が北へと伝播した。ムハンマドの時代以前に既に知られていた場合でも、その証拠は現在失われている。しかし、第八世紀にはバグダードまで持ち込まれたことは確実である。その一世紀の初めには、ムハンマド教徒が北部インドに初めて足を踏み入れ、その後数百年間、様々な運命を経ながら支配期を維持した。彼らはコルサーンやアフガニスタンも征服し、インドの学術と商業習慣が直ちにメソポタミアや西アジアの新設された学校やバザールへ容易に広がった。その道筋はキーベル峠経由でカブール、ハラート、コルサーンを通じるか、オルマズ海峡を経て波斯湾入口のバスラへ海路を通じてバグダードへと向かうものだった。実際には、現在アラブ人によって使用されている一形式のアルファベット数字符号はカブールの影響によるものであり、もう一つ最終的に私たちの数字符号となったものは別のルートでアラビアに伝わった可能性が高い。特に我々の関心を引くのは、第八世紀後半にアル・マンサールによって古い町の跡地に建設されたバグダード、即ち「平和の家」ダール・アル・サラムである。この地域は交易路が交差する場所であり、大きな都市を生み出す。実際、バグダードとアテネどちらもカーディナル・ニューマンは次のように言えるかもしれない。

アラビア世界においては、その近隣の神秘的な東方の伝統に接する利点があった一方で、接近性という面では欠けていた。また、その地には美しさが溢れており、理想郷のような土地として迎えられた。ここには偉大なるものと美しいものの原型が実在し、真理のすべての分野を探究し、知的な能力のすべての種類が示される場所があった。ここでは、芸術と哲学は王宮の如く壮麗に君臨し、精神のみが統治者であり、才能のみが貴族であった。教授たちは指導者であり、王侯は敬意を表していた。

この地にはアルマヌル、アルマムーン、ハールン・アルラシード（アロン・ザ・ジャスト）らが一時的に全人類の知性活動の中心とし、特に数学の分野においてその地位を確立した。『シンヒンド』がバグダッドに運ばれた後、モハンマド・イブン・ムサ・アルカホラジミ（アッカホラジミ）はその都市へ招致された。彼は当時の東西問わず最も有名な数学者であり、算術、日晷、天文学儀器、歴法、幾何学、代数学の著作を執筆し、アルキリツミという名前を通じてラテン語に翻訳され、新しいヒンディー数字を使用する初期の算術に「アルゴリズム」という名称を与えた。

彼は直ちに最近インドから持ち込まれた位置表記法の価値を理解し、これらの数字符号に基づく算術著作を執筆した。この著作はアデルハルト・オブ・バース（1180年頃）の時代にラテン語に翻訳され、おそらくその同時代の同国人ロバート・チェストレシスによって行われた可能性がある。この翻訳はケンブリッジで発見され、ボンコンパギーニにより1857年に出版された。

アラビア世界において、アル＝カフワルズミと同時期に活動し、アル＝マムンの下でも働いていたユダヤ人天文学者アブ・ル＝テイイブ・センディーン・イbn アリは、カリフの要請によりイスラム教に改宗したとされる。また、ヒンドゥー算術に関する著作を著しており、当時この分野に対する関心が高かったことがうかがえる。実際、ヒンドゥー数字符号をアラビア語に置き換える努力は長く続き、アル＝ナサウイの約1025年頃の算術においても、数字符号の表現法がまだ定着していなかったことが述べられている。つまり、アラビアでは二種類の数字符号が存在した：現在使用されている形式とアル＝カフワルズミが用いた形式である。

この第二形式はどのようにヨーロッパに伝播したのかという問題は残され、次章で検討される。

*       *       *       *       *

\clearpage

\chapter{第7章：ヨーロッパへの数字の本格的導入}
\section*{THE DEFINITE INTRODUCTION OF THE NUMERALS INTO EUROPE}

第七章

ヨーロッパへの数字の本格的導入

ボエティウスがヒンディー語の数字形を知っていたかどうかは確実ではなく、またゼロが含まれていなければ当然であるから、これらの数字がヨーロッパに本格的に導入された時期について検討する必要がある。東と西洋との貿易関係、そしてこれらの数字が原産地から各地の商業中心地へ移動したのは商人ではなく学問の士であったという可能性を考えると、その最初の不顕の登場時期を知ることは決して容易ではないことが明白である。今日ニューヨーク、ロンドン、ハンブルク、ロテルダムなどの商港で海員や商人が語る東洋や熱帯からの興味深い習慣について、ゲーム、社会的な特異性、衣服の奇特性などを含むような慣習は、まだ印刷物として学問の士によって記述されていないものであり、数十年、あるいは永遠にその存在が明らかにならない可能性もある。そしてこのような状況が現在であるならば、10世紀前の活版印刷発明以前には学問が最も衰微していた時代においては、その真偽を判断する上ではさらに困難であったことだろう。この文化の評価が最も低い時期に、ヒンディー語の数字はヨーロッパで初めて現れたものである。

モーレスは最初、王位主張者の補佐としてスペインに進出し、その後征服者として残った。ゴートの勢力は、3世紀間スペインを支配していたが、ジェレシー・デ・ラ・フロンテーラの戦いで10世紀初頭に壊滅した。その直後、モーレスは500年間にわたりスペインを、さらにグラナダではもっと長い期間支配し続けた。

850年までの間、キリスト教徒は宗教と政治的職位について完全な自由を享受していたため、司祭や僧侶もしばしばラテン語とアラビア語の両方を操る者となり、公式の翻訳者として活動し、自然にローマへ直接または間接的に報告を行った。実際には、この状況は永久的なものではなかったが、その期間中、東洋の学識と習慣はキリスト教スペインにもたらされた可能性が高い。

この時期には、ゴバラ数字符号がおそらくその国に存在していたと考えられ、これらの数字符号はコルダバ、グラナダ、トーレドの学校からヨーロッパへと伝播したかもしれない。

さらに、東方の数字符号が旅行者や大使の旅を通じてヨーロッパに伝播する機会は十分にありました。そのような物語を伝える記録は、 ninth century 代の有名な Arab 貿易人 Sulaiman the Merchant のものから一部を取り入れたものです。このような商人は、出会った人々の数字符号について特に詳しい可能性が高いでしょう。彼のような人物がヨーロッパの市場にそのような象徴を運んだかもしれません。

やや後には、Bagdad に住む Abū 'l-Hasan 'Alī al-Mas'udī (没年 956) が東へ China Sea まで、少なくとも Zanzibar まで南下し、西へは Atlantic まで旅をしました。彼はインド人が計算に用いた九つの数字について言及しています。

バグダードの商人アブ・アルクサイム＝ウブアイドallah アル＝イbn アル＝アハメドは、そのペルシア語名であるイbn キルダベヒによって知られるが、西暦850年頃に『道と地域の書』を著し、次のような図説を含んでいた。その中で、「ユダヤ商人はペルシャ語、ローマ字（ギリシャ語およびラテン語）、アラビア語、フランス語、スペイン語、スラブ語を話す。彼らは西から東へ、また東から西へと移動し、陸路や海路で往復する。彼らはフランスの西部海上に船を乗り継ぎ、ファラマ（古代ペルリウシウムの遺跡近く）まで航海し、そこで荷物を商隊に移し、紅海沿岸のコルゾム経由で陸路で進む。その後、東洋の海上に再び乗船し、ヘジャズとジダへ向かい、インドや中国へと到着する。帰還時には、東洋産品を運んでくる……」このような旅はまた陸路でも行われる。商人たちはフランスやスペインを出発し、タンジェに渡り、アフリカの地域とエジプトを通ってラムレールへと進み、ダマスコ、クファ、バグダッド、バスラを訪問し、アフワーズ、パルス、カルマン、インドシナを penetrare し、インドや中国に達する。このような旅行者たちは西暦900年頃には、価格の記録や市場計算において使用されるすべての数の体系について広く知るようにならざるを得なかった。この移動に関する興味深い証拠は、大英博物館に収蔵されている十字型のブローチである。これはおそらく11世紀以前からイギリスのものであり、その表面にはペルシック文字で「神以外には神はない」というムハンマディー教の刻印が施されている。アルクイン・オブ・ヨークの時代にこのような刻印がイギリスにどのように伝わったか？そしてこれらのペルシック文字が伝わったなら、なぜ数字符号もまた同様に伝わらなかったのか？

甚至高等文学作品中，亦时有迹象表明，通往远东的伟大商路从未长期关闭，贸易习惯与标记代代相传。ペルシャの詩人サディの『グリストアン』には、以下のような文節がある：

「私は一人百四十一頭の駱駝と五十人の奴隷と運送夫を所有する商人に出会った……彼は私に対して答えた：「パースの硫黄を中国へ運び、その国で高い価値があるというので、そこで中国製品をロムへ持ち込み、そこからインドへの絹織物を積み込むつもりだ。そしてハリブを通じてインド鋼（プラバ）をハーリブへと販売し、その後ヤマンのガラスを運び、ヤマンの染色布をパースに持ち帰る計画だ。」」［404］反面、これらの商人は知識人ではなく、そのような知識を論説の中で保存することも稀で、学問の人々にはたまに興味深い情報として耳にすることがあった。

東と西との間には時折使者が行き来しており、特にこれらの数字符号がヨーロッパに導入される時期にはその頻度が高まった。カールマグヌス（約800年頃）は、数学活動が始まる直前にバグダードへ使者を派遣した。そして、このような使者と共に冒険的な学問の探求者たちが同行し、アルクインがヨーク大主教アルベルト（766-780年）について述べているように、他国の知識を求めた。

さらに、この時期に東アジアとインドで設立されたニストロニア共同体はペルシャ人およびムハンマドの征服者たちから優遇を受けた。シリアのニストロニア patriarch タミオテウス（778-820年）はインドと中国へ宣教師を送り、後者の地域には主教が任命された。9世紀に中国を訪れたイブン・ワハブは、皇帝の宮廷でキリストや使徒の像を見つけたという記録がある。このような知識豊かな人々は彼らが勤務する地域について深く理解しており、帰国した際に異なる習慣も伝えずにはいられなかった。

また、アルフレッドの時代（849-901年）にはイギリスからインドまで使者が派遣され、中世を通じてアジアからの香料や織物がヨーロッパに流入し、現在では植民地やアメリカからも同様の品々が供給されるようになった。シリア、小アジア、キプロスは砂糖と羊毛を提供し、インドは香水やスパイスを生産した一方で、豪華な絨毯はペルシャと中国からヨーロッパへと運ばれた。

ジュスティニアヌスの時代（約550年頃）には既にシルク貿易が存在し、その国はシンハラと商取引を行い、さらにトルキスタンまで他の商人を通じて東方製品を西へと運んだ。7世紀にはペルシャとインド、またペルシャとコン斯坦丁堡との間で明確な貿易が存在した。ビザンティウムの $commerciarii$ は要塞に駐在し、単なる税関官吏ではなく政府の購買代理者として機能していた。

第十章においては、これらの貿易路を通じて実学に通じた商人が往来し、そのような人物は多くの習慣の知識を両方方向に持ち帰る可能性が高い。十一世紀初頭にイタリアで数字符号の一部が理解されている時期には、アフリカとアジア、さらにはインドまで旅したコンスタンティヌスというアフリカ人の人物が存在した。彼は三十九年間をかけて旅行し、バビロンでは優遇を受け、帰国後サレルノにおいて大いに歓迎された。

また、第十世紀におけるこの交流の興味深い例として、オットー一世の息子（936-973）がコンスタンティンポリス出身の王女と結婚したことが挙げられる。この君主はスペインのモーセとの関係を保ち、国外から多くの学術者を招き入れた。彼の東西との交流は、それぞれの学問を融合させるのに寄与したと考えられる。

コンスタンティヌスの洗礼（紀元312年頃）直前の時期から、数学や哲学の普及に加え、一定の文化の伝播にも寄与した有力な手段が、キリスト教徒の巡礼旅行として現れ始めた。この特徴は初期キリスト教の熱心者に見られ、その前年のユダヤ人先祖においては耶路撒冷への年一度の巡礼において僅かに見られ、他の前キリスト教の民族では殆んど見られなかった。これらの早期の巡礼者の代表として、プラチェンティア出身の二人、ジョアンとアンтонイヌス（紀元303年頃）を挙げることができる。彼らは耶路撒冷への旅の中で、後数世紀後に十字軍に至る運動を開始したように思われる。333年にボルドーからエルサレムへの最初のキリスト教徒向けガイドブック『ボルドーからエルサレムへの行程』がボルドーの巡礼者によって編纂され、その以降、聖なる巡礼は全く中断することなく続けられた。

CHRISTIAN 欧洲に数字符号が導入される確実な経路の一つは、アラブ人が地中海北部岸で行っていた掠奪と取引によるものであった。652年、ハジラル30年に至るや、ムハンマドアン軍はシチリア島の海岸に上陸し、多くの財宝を没収した。不幸にしてコンスタンスが悲惨な運命を遂げて若きコンスタントゥス4世が君臨するや否や、彼らは再び島を攻撃し、古代シラクサを掠奪した。827年、アサドの率いる下で海岸線を荒らした後も、 Syracuse を陥落させるには至らなかったものの、彼らは直ちに島の大半を掌握し、やがてその城塞を包囲するまでに成功した。しかし Syracuse が陥落する前に、彼らはイタリアの海岸線を掠奪し、ローマの壁までも到達していた。849年にレオ4世が放棄された防衛施設を修復しなければ、その年におけるイスラム軍の襲撃の影響は甚大なものになっていたかもしれない。

この時代にバグダードを去ったイブン・アル＝ハルディーバは、サラセンのパレルモで当時の大きな商業活動を描いた。同様に、9世紀後半にはピエモンテにも彼らが定着し、906年にはトゥリオを掠奪した。ソレンティノ半島では、山を登って美しいラヴェッロを見渡す旅人でも、アラブ建築の痕跡を見つけることができる。900年前後、アマルフィはモーゲルの商業中心地であり、この時期だけでなく、1世紀前から北インドの芸術家たちはこのような中心地で作品を販売していた。ハルン・アル＝ラシードとチャールマーニュの宮廷でも、彼らの作品が展示されていた。

アラブ人は、ヴェネツィアが東方の豪華さを掌握し、西方の守護者となった以前から地中海を支配していただけでなく、

「東方の豪華さを掌握し
 西方の守護者となり」

そしてジェノアがその強力な競争相手となる以前から、すでにそうした状況は存在していた。

ニコロとマフェオ・ポーロの兄弟はやや後になって、有名な旅に出発した。1260年、コンスタンティノープルを出發し、アゾフ海を経由してボカールへ向かい、その後元朝の忽必烈 Khan の宮廷に到着し、中国を penetreted し、1269年にアクレ経由で帰国した。この旅は二年後の中国への再訪問のために委任状を持参したものであった。その際にはニコロの息子マルコ・ポーロが歴史家として同行し、パミール高原を横断した。彼らは約二十数年の間中国に滞在し、海路でペルシャ経由で帰国した。

ポーロ兄弟の冒険は長く独占的ではなかった：13世紀末にはローマの宣教師と商人格レゴリウス・デ・ルコロンゴが中国に penetratet していた。1350年以前には、宣教師団体は大きくなり、信者も多数いた。東方でキリスト教徒商人集団が組織され、イタリアの中国への貿易と旅行は商業指南書の二章を占めていた。

したがって、中世においてボエティウスの時代でも、問うる学術者が商取引のために使用されるような東洋の数字符号について知るためには容易なことであった。そしてゲルバートが最も単純なこれらの形式しか理解せず、その完全な意味を把握していなかったという事実は、彼がこのように得たものであることを証明している。

ゲルバートがスペインから東方の数字符号の知識を持ち帰らなかったとしても、彼が商人物品上のマークに傾倒すれば、容易にその知識を得られたかもしれない。このような知識はイタリアの各地でおそらく入手可能であったが、それが単なる商業的な知識の一端であっても、形式は学術的な筆記体によって保存されるまでにはすぐに失われてしまう可能性があった。この時代における貿易は停滞していなかった。十一世紀と十二世紀にかけて、スラヴ人は非常に大きな商業的利益を享受し、その取引はキエフやノヴゴロドに達し、そこから東へと延びていた。コンスタンティノープルは東方との貿易の中心地であり、ビザンチン商人たちは東方の人々の様々な数字符号について完全に馴染みがあったことだろう。十一世紀にはイタリアのアマルフィがコンスタンティノープルに工場を設け、アンティオキアやエジプトとの貿易関係を持っていた。九世紀頃からヴェネツィアはシリアとカイロとの間に貴重な貿易を行っていた。ゲルバートが亡くなった後五十年、キャンターウッドのデンマーク人やノルウェー人が北 Seas に遠くまで商売を拡大し、ロシアを通じて東方へと caravan を派遣すると共に、大胆な船で直布罗陀海峡を通過して地中海へと進出していた。それからわずか数十年後、おそらく十二世紀前半に、トマス・ア・ベケットの侍従として彼の死に立ち会った書記が殉道者の生涯についての伝記を書き、その中で（幸いにも私たちの目的のために）ロンドン市への短い称賛文を付した。この書記はウィリアム・フィッツ・ステーファンと名乗っており、彼は如き言葉でブリティッシュの首都について述べている：
「ロンドンは、その繁栄と活気により、他のどの都市よりも優れている。」

アラブは金を送らず、システィナ人は武器を提供せず、エジプトには turquoiseのみの宝石が存在したとしても、中国（Seres）は紫色の衣服を輸入し、北欧はサバンナや他の毛皮、フランスはワインを輸入するなど、当時の国際貿易は広範囲にわたっていた。

その後、十字軍が東西を結びつける役割を果たした。東洋の精神は troubadour の歌において現れ、バードカイン（Bagdadの canopy）はイタリアの教会で一般的となった。シチリアとベネチアでは東洋の織物業が発展し、ス堪的ナビア半島にまで進出し始めた。

したがって、この状況は次のようになる：ヒンドゥー数字符号がキリスト教ヨーロッパの手稿で初めて現れる前数世紀には、東西間の交流が頻繁に行われていた。少なくとも9世紀頃からイタリアのような国では多くの商人にこれらの数字符号は知られていたことになるし、おそらくそれ以前にも知られていたかもしれない。しかし、これらを論理的に記述する理由はなかった。したがって、ゲルバートのような人物が学術的な集団にこれらの知識を広めたとき、彼はただ、異なる職業分野で多くの人々にとって馴染みのあるものであることを説明したに過ぎなかった。

ゲルバートが長きにわたりイタリアに数字符号を導入したとされており、彼の特異な生涯について簡単に述べるのも適切であろう。貧しい家庭に生まれたこの卓越した人物は、王の顧問となり、ついにはシルベスタ二世として1003年まで教皇冠を被り続けた。彼は幼少期からオーリラックの僧侶たちの影響を受け、特にオド・クランジーの弟子であったレイムンドのもとで学び、その結果、自己も聖職に就いた。967年頃、ボレル公と共にスペインを訪問し、ビッチのハト主教の下で三年間学んだ。ビッチはバルセロナ県の都市であり、当時完全にキリスト教徒が支配していたところであった。ゲルバート自身の証言によれば、彼の教育にはキリスト教文明の影響が大きく及んでいた。したがって彼はしばしばボエティウスについて語り、もしボエティウスが数字符号を知っていたなら、ゲルバートもその知識を得たであろうと述べている。ゲルバートがモロッコの学校で学んだ場合、ムール君主国ヒシャム王（787-822年）の令によってアラビア語を習得しなければならず、これには彼の三年間のスペイン滞在の大半が費やされることになる。しかしながら、そのような痕跡は彼の手紙には一切見られない。一方で、バルセロナは当時唯一のキリスト教徒支配地域であり、モロッコ文明と直接接触していた。さらに、同じ世紀初頭にアストゥリアス王アルンソ（没910年）は、サラゴサのワリの宮廷におけるアラブ学者たちに息子オドノの教育を委ねており、キリスト教徒とモロッコ人との間には多少の友好関係があったことが知られている。

ゲルバートはスペイン滞在を三年にわたって経過した後、紀元970年頃イタリアへと移り、ジョアン十三世大司教と出会う。彼はオットー一世皇帝の前で紹介され、翌二年（紀元972年）にルーヴィンへと招かれて哲学を学び、その地を教育中心地とするために助言した。そして1083年にボビオの住職となり、翌年ルーヴィンに戻り、紀元991年にアリウス大司教となる。政治的理由により紀元996年にイタリアへと再び移り、紀元998年にラヴェナの大司教となり、翌年には教皇に選出された。

彼は遠く未来の智者でありながら、そのような学者たちが近世以前にもしばしば遭遇したように、異端罪や魔術の疑惑で苦しんだ。1522年のヴェネツィア発行の生誕記録によれば、黒魔法によって教皇に昇進したとまで言われている。しかしゲルバートは占星術にも関心があったが、これは当時の天文学であり、学問的な人間であれば知る価値のある科学の一つであった。

ゲルバートとその弟子たちはゴバラ数字符号を知っていたという事実は、今や論争の余地がないものである。ベルネリヌスとリッチャーは彼らを「キャラターゼ」という有名な名称で呼んだが、ラオンのラドウフもその意味での「キャラターゼ」を使用したのは百年後のことである。おそらくゲルバートはこれらのゴバラ数字符号をクリスチャンヨーロッパにおいて初めて科学的な方法で説明した人物であり、ゼロを含まなかった。もしゼロについて知っていたとしても、その使用法には理解していなかったと考えられる。

彼がこれらの数字符号を得た場所については、まだ確定していません。
スペインから持ち込んだわけではなかったことは、多くの慎重な調査家たちの見解です。これは、ゲルバートが在籍した後期に学問で有名となった人物たちがほとんどいたためと考えられています。そのような人物たちは、初期にはイブン・サイナ（アビシーナ）、中期にはセビリアのゲルベル、そして最終的には十二世紀末期にアブ・ロシュド（アバローゼス）です。一方で、彼が三年間近隣のアラブ人と接触していたことから、彼らの学問を吸収した可能性があるという見解もあります。当時のキリスト教徒とモーレンとの境界は明確に区別されていました。

しかし、商業的な数字符号システムは商人によってより容易に知られるべきものであり、学者たちよりもそのような知識を得る可能性が高いと考えるべきです。巡回販売人はスペインでも禁断の領域を知らず、他の地域でも同様でした。ゴバール（.gobār）の数字符号が商品を記録するか単純な計算のために使用されていた場合、彼はそれらを知っていたでしょうし、若きフランスの修道士の探求心にこれらの知識を伝える人物としては彼の方がより適切でした。ゲルバートがこれらを不完全にしか理解していなかったこと、そして彼がそれらを使用したのは計算のみであり、またその形態は明らかにスペインのゴバールと似ていることから、これはモーレンの小さな商人を通じてこの多才な学者が知識を得た可能性が高いと言えます。

さらに、彼名で冠された幾何学の部分があり、それが彼のものであることが疑いようがないものも存在し、アラブの影響を示しており、少なくとも彼は移植された東洋の学問と接触したことを証明しています。完全には理解していなかったとしてもです。当時から現在まで、両方の人々と交易するPersistentなユダヤ商人がおり、有用な知識を得ることに常に熱心でした。このような源から学問を喜ぶ人物としてゲルバートは自然でした。

一方、ゲルバートの生涯に関する主要な情報源は、彼がスペイン滞在中にモーレーゼの影響下にあったことを示す具体的な証拠を提供していない。これらの情報源は彼の手紙とリシャールによって書かれた歴史である。ゲルバートは書簡術の達人であり、その尊い地位により、彼の手紙がその優れた古典性にかかわらず保存された。リシャールはレーニュ聖ラミ教会院の僧侶で、おそらくゲルバートの弟子であった。ゲルバートがレーニュ総主教となった後、彼から歴史を書くよう依頼され、その要求に応じた。この作品は手稿として保存され、1833年にペルツがバムブルクで発見されるまで完全に忘れ去られていた。この著作はレーニュ総主教のゲルバートに捧げられており、彼がモーレーゼの学問を獲得するために行った努力について証言していたこと自体が確実であったと思われる。

現在の事実は、ゲルバートがサラセン人との接触について述べていないことである。彼の手紙にはモーレスやアラブ数字、コルドバに関する言及がない。スペインはその名前で呼ばれておらず、スペイン人の学術者も僅かにしか登場していない。彼の手紙の中でジョセフ・イスパヌスまたはジョセフ・サピエンスについて述べているが、この賢者のジョセフとは誰であるかは不明である。おそらく彼が若きフランスの僧侶に不完全な知識を提供した人物であった可能性がある。ゲルバート訪問後数年以内に、名前には劣る二人の若きスペインの僧侶が、数学的な興味は彼より薄かったとしても、彼らが持っていたヒンディー数字形についての僅かな知識を手稿として永続的に記録することに価値があると判断した。そのような知識が北スペインの簡素な修道院アルベルダやアルバーダにまで浸透していたことは、それなりに広く普及していたことを示している。

ゲルバートの著作『数の分割に関する小論』はチャスレによって「科学史の中で最も不明瞭な文書の一つ」と評価されている。ゲルバートの他の数学的著作について最完全の情報は、ブーブノフによって提供されており、彼はこの作品を真実であると考慮している。

ゲルバートはこれらの数字符号をいかに重んじていなかったか、彼の『Regula de abaco computi』や『Libellus』といった著作において、これらを使用することなく、ローマ字の形式のみを使用している。しかし、ベルネリヌスは九つの「ゴバール」文字について言及しており、ゲルバートは千枚の「ジェトン」またはカウンターやアバクスにこれらの文字をマークしていた。彼は占星術師にアバクスを作成させたが、十の列には八個のカウンターではなく、単一のマークされた8のカウンターを置き、他の場所も同様にした。ゲルバートはゼロがないため、必要とする場所にはカウンターを配置できず、その空白をゼロで埋める方法がなかった。

彼はおそらくこれらの「キャラターゼ」を呼んだと考えられ、この名称は数字符号にも適用された。『ブエティウス』の補遺において「アピックス」と呼ばれるこれらは、一般的に銭貨学において「テッサレイア」として知られるローマのジェトンと関連があるかどうか興味深い推測である。これらのジェトンにはI-XVIまでの数字が刻まれており、十六はセストルティウスにおける-Assi-の数を指すものと考えられている。

『アピックス』上の数字にはイギン、アンドラス、オルミス、アルバラス、キマス、カルクトスまたはカルティス、ゼニス、テメニアス、セルエンティス、シポスといった名称が与えられ、これらの語源と意味はまだ謎に包まれている。いくつかの単語にはシトミック起源があると考えられる。ウハド、タンイネ、タルタ、アルバ、クムサ、セッタ、セブバ、ティミニア、タセードは、モロッコで使用されるアラビア方言における一から九までの数字符号の名称とされ、その中には四、五、八の単語が特に類似している。

名前「apices」は後世では一般的ではなく、代わりに「notae」という言葉がしばしば用いられ、最終的にはこの言葉から記号法の名称が生まれた。また、「figures」「ciphers」「signs」「elements」「characters」といった言葉も頻繁に使用された。

ゲルバートの教えによって新しい数字符号を広めることはほとんどなかったため、オーキレット（O'Creat）は1世紀後輩でありベス Thatcher の友人兼弟子であるアデハールドと異なり、ローマ字とゼロを使用した。ゲルバートはゼロを使わずゴブラー形式の数字符号を用いていた。

オーキレットはゼロに3つの形を用いた、円点（o）、丸（＝o）、ギリシャ文字のτ（[Greek: t]）。この使用法には当然、位置価値が伴い、彼はIII III を 33、ICC OO と I. II. τ τ を 1200、I. O. VIII. IX を 1089、そして I. IIII. IIII. ττττ を 1200 の平方数として記した。

ゲルバートの時代からレオナルドの巨著が出版されるまでの期間は「算術家」の時代と呼ぶことができる。十二世紀初頭にヒンディーの計算法を説明する書籍が出版された後も、長い年月間、算術家、骨算盤を主張する者たちとアルゴリズム派、新しい数字符号を支持する者の間で論争が続いた。「cifra」と「algorismus cifra」は一種の嘲弄的な意味を持ち、実質的にゼロは骨算盤においてどの位かを示す列によって与えられる任意の単位の価値を決定するため無用であることを示唆していた。

このように、ガウチール・ド・コインシー（1177-1236）はマリアの奇跡についての著作で以下のように述べている：

「角のある獣、羊
  アルゴリズムス・シフラ
  祭日において聖母を祝わない司祭」

算盤は長きにわたり、新しいアルゴリズムを用いた新数字符号に対してもその地位を譲らなかった。チャーチャーの『ミラーの物語』においては、彼が次のように描写している：

「彼のアルマゲストと大本小本の書籍、
  彼の星象盤、芸術のために欲するもの、
  芸術石を整然と並べた棚、
  床頭に置かれた」

またチャーチャーは、彼の息子レオワのために書いた星象盤に関する説明においても、儀器上の度数はヒンドゥー・アラビア数字で表されている：「その上には五分ごとに区切られた芸術石の数が刻まれている」、「これらの数は芸術の方法で記述されている」という文句があり、これはアルゴリズムの方式を指している。1387年頃、トーマス・ウスクは次のように述べている：「自身には意味を持たないシフターも、他の数字符号に意味を与える力を持つ」。新しいアイデアの吸収はいかに遅く、いかに苦痛であるかを示すものである。

ベルニリヌスは、算板が WELL-POLISHED BOARD（精緻な板）であると述べており、その上に BLUE SAND（青砂）を敷き詰め、幾何学者が幾何学的図形を描くために使用していることを記述している。我々は以前にも述べた通り、ヒンズーも計算のために砂で計算を行っていたが、彼らが列算板を持っていたという具体的な証拠は存在しない。

ベルニリヌスによれば、計算のためには板は三十三の縦列に分割され、そのうち三列は分数用として予約されている。単位列から始め、Bernelinus は各三列を SEMICIRCULAR ARC（半円弧）で区切るという言葉を使用している。さらに、単位列には小さな半円弧が置かれ、十と百の列には別の半円弧が接続されている。このようにして、PICTAGOREAN ARCUS（ピタゴレアの半円弧）または単に ARCUS という名称が生まれた。

この形式の算板上で加法、減法、乗法を行う操作は説明を要さなかったものの、特に異なる位数の乗法については詳しく扱われていた。しかし、除法の操作は困難であった。補余差による除法の方法についての説明は、中世の算術学者にとって長年障害となっていたが、その詳細については数学史に関する著作および算板に関する特別な著作を参照されたい。

主題に取り組んだ著者には、フリューのアボ（970年頃）、ローブスまたはラウバッハのエルジェリク（950-1007年頃）、そしてコントラクトゥスのエルマヌス（1013-1054年）が含まれる。これらの人々はすべてローマ数字のみを使用した。同様に、バースのアデルハルト（1130年頃）は『算盤規則』において新しい数字を参照していないが、彼はそのことを知っていたことは確実である。この時期にヒンドゥー数字の一部を使用した計算板に関する他の著者には、12世紀前半のゲルランド（491）と1200年頃のターチール（492）が含まれる。この時期に使用された数字の形については、88ページの図を参照されたい。

注：文中の「アボ」、「エルジェリク」、「エルマヌス」、「ゲルランド」、「ターチール」は既存の定訳を使用していますが、原文中のカナ表記をそのまま使用しました。正式な定訳が確認できない場合は、適切な翻訳やカナ表記に修正することが推奨されます。

ゲルバート没後、少しずつヨーロッパの学術者たちは新しい数字を知り、主にアラブ学問の導入を通じてその認識が深まった。中世の暗黙期は過ぎ去ったものの、数学者としてのゲルバートのような顕著な支持者は二百年間にわたって現れなかった。この偉大なる復興について、十一世紀のクルニ大聖堂の Benedictine 修道院所属のラウール・グラベル（1085年頃 - 約1046年）は次のように述べている。「世界が古い時代の着古れた衣服を捨て去り、教会の白い衣装に身を包もうとするかのようだった」。この宗教活動とともに、他の分野への関心も高まった。アルゴリズムが現れ始め、外部からの知識は興味深い聴衆を得た。ラウール（またはラドゥルフ）と名乗る者、彼の故郷の修道院学校で教師を務めた者であり、ラオヌスの安セム（有名な神学者）の兄弟である者によって、音楽に関する論文と算術が書かれ、ナグルは1890年に前者を発表したが、後者は七十七葉の羊皮紙手稿に保存され、ローマ数字と [.g]ob[=a]r 数字の奇妙な混合体を含んでおり、前者は大きな結果を表現するために使用され、後者は実用的な計算のために用いられた。これらの [.g]ob[=a]r 「キャラクター」にはゼロ（sipos）が含まれており、ラドゥルフはその完全な意味を知らなかったことを示し、十二世紀の初頭において中央フランスの教会学校ではまだ体系化されていなかったことが明らかである。

この時期に、アルゴリズムとシフраという言葉は数学的文献以外でも一般的に使用されるようになった。ジョーダンはアラヌス・オブ・インソイルス（アラン・ド・リール）、ガウティエ・ド・コincy（1177-1236）などの著作からそのような使用例を多数引用している。

この興味深い時期には、ラーナ、シベリュのジョアン、ヒスパレニカスのヨハネ、トレタナウスのヨハネ、そしてラーナのヒスパレンシカスと呼ばれたスペインの優秀なユダヤ人ジョアンも算術に貢献した。彼の生年は、ローテルダム総主教ライムンドへの作品献呈という事実により、1130年から1150年の間に存在したことをほぼ確実にしている。彼の関心は主にアラビア語著作の翻訳にあり、特にアリストテレス哲学に関連するものが多く含まれていた。

算術の観点からは、ジョアン・ヒスパレニカスの『アルゴリズミの実践的算術』という題名の手稿がボンコンパギーニによってパリの現在の国立図書館で発見されたことにより、主な興味が集まっている。この著作はアルキワルズミの作品であると明確に主張しているが、その証拠は一貫して否定的であり、しかし本自身は非常に価値がある。なぜなら、それは著書が記された時代における知識を反映しているからである。この手稿は整数および六十進分数（包括的に根号も含む）の演算について述べており、実用的な応用例は含まれていない。

ラーナと同時期にトレタムナで生活していたゲハルト・デ・クリモーネもまた、彼はゲナルドゥスと誤って同一視されることもあるが、これは誤りである。彼は医者であり、天文学者でもあり、数学者でもあり、イタリアとスペインでアラビア語から翻訳を行っていた。算術において彼はアルゴリズムの思想を広める上で重要な役割を果たした。

十二世紀にドイツで最初のヒンディー語数字符号を使用した計算の痕跡は、ウィーンの宮廷図書館に収蔵されている1143年頃の『アルゴリズムス』に見ることができる。この文献は同じ著者の『コンピュトス』と同封されており、その日付が記されている。本書には「加法について」、「減法について」、「平均法について」、「除法について」という章があり、また乗法に関する一部の章も含まれている。

数字符号は通常の中世の形態をとっているが、2は図表から見て異なる形をしており、3は特徴的な形状 [符号] を取っている。この形状は十二世紀に見られる特有のものである。

平安時代後期から鎌倉時代にかけて、数の書籍『数書』がヘブライ語で出版された。著者はラビ・アブラハム・イブン・メイル・イブン・エズラであり、1092年頃にトーレドに生まれた。1139年にエジプト、パレスチナ、東方を訪問し、イタリアでも数年の間滞在した。その後、南部フランスとイングランドで生活し、1167年に没した。彼がヒンドゥー算術の知識を得た可能性は、故郷であるトーレドで学んだものであるが、旅行中に得た他のシステムの知識も、この算術に対する理解を深めた可能性がある。

彼は、希伯来語アルファベットの最初の文字を用いて数を表し、[ヘブライ: A B G D H W Z CH T`]で9 8 7 6 5 4 3 2 1に対応させた。また、ゼロは円で表した。付記に引用された文書によれば、彼はヒンドゥーの起源を知っていたが、手稿ではヒンドゥーの形式を記載しながらも、上記の九つの希伯来文字と位置価値を使用して計算を行った。

\clearpage

\chapter{第8章：ヨーロッパ各地における数字の普及と計算革命}
\section*{THE SPREAD OF THE NUMERALS IN EUROPE}

第八章

ヨーロッパにおける数字の普及

中世の著者の中で、おそらく最も新しい数字符号を欧州の学術界に導入した人物はピサのレオナルド・フィボナッチである。この驚くべき人物は、中世最大の数学的天才であり、約1175年頃にピサで生まれた。

当今の旅行者はフィボナッチ通りを通過し、サンタ・コルソへ向かう途中でラーディノ数学者の像を見ることができる。この町は数学家中に最も崇敬の念を表していることから、レオナルドはこの都市の黄金時代に生まれた人物であり、その時代は商業、宗教、知性の繁栄期であった。

ピサはアーノ川の口に位置し、十三世紀初頭にはジェノとベネチアと共にイタリア最大の商都 trio を形成していた。ベネチアがレバント貿易を掌握する以前から、ピサは東方との密接な関係があった。古いラテン歴史書によれば、1005年には「ピサはサランクスに占領された」と記されている一方で、翌年には「ピサ人はレッジオでサランクスを破った」、そして1012年に「サランクスがピサを攻撃し、その町を破壊した」という記述がある。しかし、この都市はすぐに回復し、1099年にはシリアに百二十隻の船を派遣し、数年後にはコンスタンティノープルに商取引所を設立した。また、イタリア内での都市間戦争がその活動を刺激していた。1114年の著者は当時、ピサには多くの異教徒—トルコ人、リビア人、パルティア人、カルデー人—が存在したと記している。このような戦乱の中で、多国籍都市であり商業中心地であるピサにおいて、文学者の作品は高く評価されなかった。そのような中でレオナルドの才能が歴史の一驚きとなり、また「占星術を描くな、期待しない方が良い、期待したことは実現しない」という警告を与えることとなった。

レオナルドの父はウィリアムであると記されているが、彼の兄弟ボナッチングスも存在する。しかしながら、レオナルドの家族についてはさらに詳しいことは知られていない。フェルナリーゾについて多くの著者は、彼の父の名前はボナチオと推測しており、これから「ボナチオの息子」またはFibonacciと呼ばれたものと考えられている。しかし、別の見解では、この名前は単なる一般的な名前であり、ボナチオというラテン語でも同様に「ボナチオの息子」という意味で使用されただけであるとし、そのような名前の唯一の当時の証拠は存在しない。レオナルドが用いた「ビゴロ」の名前については、自身が付けたものであり、頭が悪かったという意味で使われた言葉で、商業界や大学団体によって彼に適用されたものと解釈され、彼はそれを取り入れて自己主張したと考えられる。しかしミランゼィは、ビゴロ（またはピゴロ）という言葉がトスカーナ地方で旅行者を指す言葉であり、レオナルドが外国で学んだ人物としては自然にその名前を選んだことを示している。

レオナルドの父は、ブギアという地名で現在のボージュに所在する商業代理店を務めており、バルビーリ海岸のサルダイ古城址として知られる古代都市でした。この地はヴァンダール朝の王都であり、レオナルドが生まれる前1世紀にはベンイ・ハムド朝の王都でもありました。ブギアは海岸線に位置し、ララ・グルア山によって守られている優れた港を持つことから、十二世紀末にはアフリカ貿易の中心地となりました。レオナルドは幼少期にこの地で過ごし、モーリシャス人の教師のもとで学びました。青年期になると地中海を巡游し、エジプト、シリア、ギリシャ、 Sicilia、プロヴァンスを訪ね、商人だけでなく学者とも交流し、貿易の中心地で用いられている様々な数の体系について知識を得ました。しかし、彼はこれらのすべての体系をヒンドゥーの数の体系と比較してほとんど誤りだと考えました。

レオナルドは1202年に『 Liber Abaci』を著し、1228年には再編纂しました。この著作では、数字が説明され、商業計算に使用されます。しかし、この著作は商人群体にとっては先進的すぎ、保守的な大学団体にとっては新奇すぎるため、人気を得ることはできませんでした。当時の数学は、パリのソロバンヌで1215年に知られる最古の条例によれば、新設された大学における位置が僅かしかなかったことからも明らかです。この時期は商業活動が盛んであり、このような著作は通常よりもさらに注目を集めることは難しかったのです。

レオナルドが知らしめた新しい数字符号がヨーロッパ西部で即座に採用されるとは考え難く、その優れた性質にもかかわらず、キリスト教徒のヨーロッパにおいては既存の数字符号よりも著しく劣っていたものと比較しても然りである。大学院の反対がこれほど改善に効果的ではないと考えるのは、二五〇年前に印刷術が発明される以前のこの時代における知識の拡散の困難さを忘れてはならない。当時、教皇や王侯貴族すらも多くの現在の農民よりも書籍の所有数が少ない状況であった。「九世紀にルッカのサン・マルティノ大聖堂教会所蔵の図書館には、神学解釈注釈からの要約本十九冊のみがあった」。実際上は、十五世紀初頭までパラ・ディ・ストロッツィがペトラーラの計画を実現するための手続を講じるまでは、公共図書館の設立が行われなかった。したがって、この早期知識は主に口承によって伝えられたのである。幸いにも当時のイタリアにおける外国留学生の存在により、その伝達が可能となった。人間の性質が今と同様であれば、レオナルドの著作に対する大学院の反対はむしろ学生たちをそれらの研究に熱心に取り組ませた可能性もある。例えば1209年ヴィチェンツァでは、ボヘミア人、ポーランド人、フランス人、ブルゴーニュ人、ドイツ人、スペイン人が存在し、イタリア各地の都市代表も含まれていた。そのような状況は他の知的中心地でも同様であった。したがって、この知識は必ず広まり、実際には多くの証拠からそれが事実であることが確認される。フローレンスの銀行家たちが1299年にこれらの数字符号を使用禁止され、パドヴァ大学の法規が「書籍の価格表を文字で明瞭に記載せよ」と求めるにもかかわらず、実際には1275年頃から数字符号は大きな進展を見せた。

ドイツの一般庶民が十六世紀以前にアラビア数字を使用したことは比較的例外的であり、これは1457年から1496年の間の流行の暦書によっても明らかである。これらの暦書は通常ローマ数字を用いており、コエベルの1518年の暦書では、アラビア数字がローマ数字に次ぐ形で使用されている。ドレスデンのクレウスシュールの登録簿においても、ローマ数字が1539年まで使用されていた。

レオナルド・ダ・ペルサの科学的業績を軽視するつもりはなくとも、アレクサンダー・デ・ヴィラ・ディ（約1240年頃）およびジョン・オブ・ハリファックス（ Sacrobosco、約1250年頃）のより一般的に使用された著作物により、一般庶民における数字の普及がより大きく促進されたことは疑いようがない。

アレクサンダー・デ・ヴィラ・ディの『アルゴリズムの歌』は詩で書かれており、そのような時代には多くの教科書も同様であった。その広範な使用は、ヨーロッパの図書館に残る多くの手稿数から明らかである。 Sacroboscoの『アルゴリズム』において、『アルゴリズムの歌』からの引用が含まれており、大学教育用教科書として広く人気があった。この著作は、その目的のために作られたものと推測され、多くの大学講師による解説が見られる。最も広く使用されたものは1291年にペトロス・デ・ダーシアの解説であり、おそらく最も一般的に使用されたものである。

これらの著作物は、十三世紀から十六世紀までの大学における数学教育方法について興味深い光を投げかける。明らかに、テキストは最初学生たちによって読み取り、抄録された。その後には、ペトロス・デ・ダーシアの解説のように、一行ずつテキストの解説が行われた。

Sacrobosコッコロスの著作は、おそらくこの著作の広範な使用により、「アラビア数字」という用語が一般的になったことに興味があります。二箇所で、このシステムの発明者について言及があります。序論においては、計算術はアルグスという哲学者によって生み出されたものであると述べられており、その名前から「アルゴリズム」という用語が派生したのです。また、数値表現の部分では、アラブ人がこの科学の発明者であると指摘されています。一部の解説家、ペトロス・デ・ダシオを含む者たちにはヒンドゥー語の起源が知られていたものの、多くの解説家は文書そのものを受け入れたと考えられるため、アラブ人がこのシステムの発明者と信用されたのです。

フランス語におけるアルゴリズムの最初の明確な痕跡は、1275年頃に書かれた手稿から見ることができます。この興味深い葉書は、アルゴリズムに関する部分が一枚の葉に収まっているため、1741年にアブベ・ルボエフによって、1824年にダヌーによって指摘されました。その後、パリの手稿群の中に埋もれてしまったかのように見えましたが、チャールズ＝ド＝シャルーズがその探求を徒労に終わらせたとしても、再発見は1882年に実現しました。その年、シャル・アンドレ・フランソワ・ド・サン＝ヒネによってこの手稿が見つかり、彼の注意深い保存のおかげで、興味深い手稿に関する知識を得ることができました。この著作は匿名であり、ほとんど完全に幾何学に専念していますが、算術に関連する部分は二ページ（一枚の葉）に過ぎません。これらの部分では、数字の形が示され、四則演算について非常に簡略な説明がなされており、著者はこの主題について僅かな理解しか持っていないことが明白です。

フランスで新しい数の体系が知られるや否や、その影響は海峡を越えてイギリスに及んだ。ヒグデンは、十四世紀初頭からそれ以降の数十年間におけるフランスの影響について述べている。「二百余年の間に、学校において子供たちは自分の言葉を使わずに、フランス語で授業を受けることになり、また自らの物事をフランス語で説明することも求められるようになった。……貴族の子弟は、胎内から始めてフランス語を学ぶことが習慣となり、上流社会の人々も自分たちを貴族とみなすために、大きな努力を払ってフランス語を習得する」というように述べている。

これらの新しい数字符号がより速やかに注目を集めるべきではなかった理由は何かと問われることがある。なぜこれらは商業や学校で一般化されるまで十六世紀を待たなければならなかったのか？これに対し、学習は常に貿易の要求に従うものであり、人々の生活に直接関連するような工作は合理的な近道を通らなければならない。十五世紀前半までの商業計算では新しい数字符号が求められなかった理由は二つある。まず第一に、安価な紙が存在しなかったことである。紙漉き技術は十二世紀までヨーロッパには導入されず、安価な紙は十九世紀の産物である。また、現代型の鉛筆も十六世紀から始まったものである。第二に、乗算や除算などの近代的な計算方法がまだ発明されていなかったことである。これらの操作はすべて複数の単位を含む尺度で行われたため、各段階で変換が必要となり、それらはより重要な作業となった。古い方法では使用した後は数字を消去しなければならなかったが、これは計算器を使用すれば比較的簡単に行うことができた。一方、新しい方法ではこの操作が容易ではなく、そのため新しい数字符号は計量所で非常に遅々として認められ、教師たちには熱意を持って歓迎されなかった。

新数法の早期著作において用いられたこと以外に、これらの数字は手稿の日付や記念碑上にも時々見られていた。ヨーロッパで最も古い確実な日付が付けられた文書として知られるのは、976年（西暦）にスペイン北部ロゴロン近郊アルバルダ修道院で書かれたラテン語の手稿であるコーデックス・ビジャランヌスである。この中にはゼロを除く九つの（ゴバール型の）文字が最初の三冊目の『原論』の章末に加えられており、ここでイシドロス・デ・セビリアはローマ数字について議論している。

同じ著作のスペイン語抄本で、992年（西暦）に見つかったものも同様の部分に数字を含んでいる。作者はこれらの文字がインド由来であると述べており、「また算術の図形について。我々はインド人が最も精密な知恵を持っていることを学び、他の民族もその算術や幾何学、以及其他自由学科において彼らに敬意を表すべきである。これがあらゆる段階のどの段階でも一つずつ示す九つの数字から明らかである」と述べている。その後に九つのゴバール文字が続く。

以前に与えられた数盤の形態も十世紀のものと疑われるものがある。新数法を含む最初のアラビア語文書は、874年および888年の二つの手稿である。これらの数字は、970年にシリアで書かれた著作で約100年後に現れる。また、最近発見されたエジプトのジェリマイヤ修道院近くに存在した柱にも、新数法の使用痕跡がある。この柱にはアラビア語の刻印があり、その日付は349 AH（西暦961年）に対応している。

ラテン語文書の日付付けにおいては、13世紀からアラビア文字が用いられた。

THESE 数字の早期使用に関するヨーロッパにおける唯一の科学的な研究は、ブリティッシュ・ミュージアムのヒル氏によってなされたものである。彼の調査によれば、これらの数字で表記された日付が刻印されているコインの最初の例は、シチリアのロジャー王在位期の1138年に見出されることがわかった。最近まで、最も古い日付は1217年（西暦）のアラビアの硬貨と1388年のトルコの硬貨であるという認識が広く存在していた。一方で、かつて非常に早期とされていたと思われていた多くの印鑑やメダルは、ヒル氏によって比較的最近の作品であることが示された。

しかし、ヨーロッパの手稿には、十二世紀以前にこれらの数字が使用されている例が多数存在する。コーデックス・ビジャヌスに見られる例以外にも、現在チューリッヒの大学図書館に収蔵されている聖ガルの手稿において、十世紀の別の例が発見された。

ヒル氏のリストにおける第三の標本は1077年のバチカン手稿からである。第四と第五の標本は、同じ十一世紀（ボエティウスの）エルランゲン手稿からであり、第六と第七の標本もチャールズ＝トリーズのボエティウス手稿（同様に十一世紀）からである。これらの早期の形態や他の例は、ヒル氏の許可を得て再現した次の表において示されている。

最初期の手稿形式

[図版]

本章はMr. Hillの貴重な論文に記載されている五十以上の表の一つであり、この単著には十世紀から十六世紀までヨーロッパにおける数字符号の発展に関する詳細が述べられている。彼は、既に述べたシチリアの例を除き、これらの数字符号で表された最初のヨーロッパのコインやメダルの年号として次のものを発見したと述べている：オーストリア、1484年；ドイツ、1489年（コリン）；スイス、1424年（サン・ガール）；ネланд、1474年；フランス、1485年；イタリア、1390年。

これらの数字符号で表された最初の英語のコインは1551年に打たれたが、スコットランドでは1539年のものがある。印刷本のページ番号を数字符号で表すのは、ペトラルカの著作『Cologne』版1471年刊行作において初めて行われた。その日付は『Biblia Pauperum』（1470年刊行）において次のように記されている：

[Illustration]

また、おそらく1430年前後に遡る別のブロック本では、数字符号が複数の図面に現れ、以下のような形態を呈している：

[Illustration]

1471年以前の多くの印刷物には手書きで番号をつけられたページや章があるが、これらの多くは著作の印刷よりも後の日付である。他の著作はおそらく印刷後に直接番号をつけられたものであった。例えば、1470年の本においては2, 3, 4, 5, 6章は次のように表されている：Capitulem [Symbol 2]m.,... [Symbol 3]m.,... 4m.,... v,... vi，そしてローマ数字が続く。これは印刷者が手書きで埋めるためにスペースを残した本体の本文に現れている。また、別の1470年の本では、ローマ数字とヒンドゥー数字の混在でページ番号をつけている：

[Illustration] 125      [Illustration] 150
  [Illustration] 147      [Illustration] 202

monumental銘文においては、カタレイン近郊のトロパウに於いて1007年とされる墓碑が存在したという説があった一方で、ビエブリヒには1299年に刻まれたものがあるとも報告された。しかし、確実なものはプフォルゼムの1371年銘文とユルムの1388年銘文であることが確認されている。ウェルズ大聖堂に於いて十三世紀に属する数字が記されたとする説もあるが、これは明らかにその後期なものである。

第143頁の表は、ヒル氏の著作における表を補足するために用いることができる。



初期の手稿形式

 [図版] 十二世紀
   [図版] 1197年
   [図版] 1275年
   [図版] 約1294年
   [図版] 約1303年
   [図版] 約1360年
   [図版] 約1442年



さらに比較のため、プラムプトン氏蔵書より再現された『 rararithmetic a』から三枚の図版を引用することができる。最初のものは、パリで1424年にロランデスというポルトガル医師によって作成され、ジョン・オブ・ランカスター公ベフォード侯に命じられた算術に関するラテン語手稿から取れるものである。数字は2の連続したべき乗を示している。第二の図版は1475年頃のルカ・ダ・フィレンツェの『 inprincipiodarte dabacho』から、第三のものは約1500年に作成された無名手稿から取れる。

[図版]

数字符号の形状については、印刷が発明されるまでには、風尚が大きな影響を与えました。これによりこれらの形状は固定されましたが、書記においてはまだ大きな変動があります。例えばフランス語の5やドイツ語の7および9を見るとその様子がわかります。印刷においても完全な一貫性は見られず、外国人にとってはしばしばフランス型の3と5を区別することが困難です。

【図版】

数字符号について、次のものは後期の手稿と初期の印刷本に見られる形状の一例です。

1. 初期の印刷本では「一」がしばしば「i」となっており、これは型を節約するために使用された可能性があります。現代のタイプライターも数字1とアルファベットの「l」を同一の文字で表す場合と同じです。手稿においては「一」は以下の形状で見られます。

【図版】

2. 「二」は初期の印刷本ではしばしば「z」となっており、12は「iz」となります。中世の手稿においては以下のような形が一般的です。

【図版】

これらの早期の痕跡から明らかに、これは原始的な形状 [二横線] の草書体であることがわかります。3も同様で、三横線から来ていると推測されます[三横線，原文為“3 horizontal strokes”]（ナナ・ガットの彫刻においても同様）。

3. 「三」は最初の印刷本では通常特別な形状を持っていましたが、まれに単なる [符号] の形で表されることがあります。中世の手稿においては他の多くのものよりも変動が少なかったようです。以下のような形が一般的です。

【図版】

四は大きく変化し、算術の手稿がどの時代に書かれ、どこで書かれたものであるかを判断する際の最初の試練の一つは、この数字符号の形態の検討である。印刷技術が普及するまでの最も一般的な形式は「」であったが、レオナルド・ダ・ペスカローニの著作のフローレンス手稿では、「」という形が見られる；しかし、手稿を調べると、フローレンスの算術学者や天文学者は比較的早くから最初のこれらの形式を直線的なものへと変更し、「」「」や「」のような形に近づけたことがわかる。これらは現在の我々が使用する形式に近い。最初の印刷物では一般的に、上端が閉じている「」という形式を使用したが、書き込み時に開いた上端（「」）を用いるものは完全に近代的なものである。以下は四の他の形態からいくつかの手稿を引用する。

[図版]

五も印刷技術が普及するまでの時期には大きく変化していた。以下はその一部の形式である。

[図版]

六は他の数字符号と比較して、あまり大きな変化を見せていない。主な変化は上端の傾斜に現れている。以下に示すようにわかる。

[図版]

七も四と同じく、15世紀以降に現在のような直立した形を取ることになった。中世には以下のようであった。



八も六と同様に、あまり大きな変化を見せていない。中世には以下のような少数の変種が存在した。

[図版]

しかし、16世紀になると、八を「」という形で書き始める傾向が現れた。

九も他の数字符号と比較して、あまり大きな変化を見せていない。中世の形式には以下のようなものがある。

[図版]

零の形状も多様な歴史を有していた。以下は一般的な中世の形態である。

[図版]

位置価の説明は、多くの初期の著者にとって重要な問題であった。彼らがロシア語の「チョトゥー」のような算盤を使用し、この算盤を位置記数法を学ぶすべての人々に提示していたら、大きな困難は生じなかっただろう。しかし、中世の縦計算では、線が10のべき乗を表し、空隙がその半分を表すため、このような比較には適さなかった。したがって、以下のような丁寧な説明を見ることができる。

「これらの数字はそれぞれ自身を表し、それ以上ではない。最初の桁に立つとき...

「第二桁に立つとき、それは十倍自身を表す。この2が20を表すように、2は二を表し、二の十倍は二十である。そして彼が左側にあり第二桁に立つとき、それは十倍自身を表す。そして続く...

「ゼロ cifra は意味を持たず、しかし次の数字をより多く表すために存在する。この詩を解説せよ。cifra は何の意味も持たないが、彼は後の数字を表現するために存在し、彼自身が欠けている場合でも、前の数字がないときには一つだけを表すべきである。例えば10の場合、この数字1は十を表し、もし cifra が欠けていて前後に数字がない場合、彼はただの一つを表すべきである。なぜなら彼は最初の桁に立つべきであるから...」

このようなシステムが文書、硬貨、記念碑の日付として使用されていたとしても、その普及は16世紀まで待たなければならなかった北イタリア以北の国々においても、より早く広く採用されるべきであったと思われる。しかし、これは多くの発明品に見られる運命であり、今日でも対数や短縮計算過程を無視しているように、そのようなことが証拠となっている。

ドイツにおいては、15世紀には新しい数字符号の隆盛が見られ、16世紀には徐々に支配力を得、17世紀にはついに2000年以上にわたって商業計算を支配したシステムを制圧することとなった。この新規画の成功の一端は、ルターの言語化要求によるものであった。すなわち、すべての学問が方言に入ることを求めていたからである。

ローマ数字からアラビア数字への移行期には、様々な異形の形式が存在した。例えば、14世紀にはc[α]で104、1000.300.80 et 4で1384を表すことが見られた。また、15世紀の手稿では12901が1291を指す場合があった。同じく15世紀にはm. cccc. 8IIが1482を表し、$M{oCCCC}$o50（1450）とMCCCCXL6（1446）はテオドリコス・ルフィーによって使用された。

次の時代には1vojjが1502を表す形が見られた。さらには、サフォルターニの『ニューボーム』においても序数の使用は混乱しており、同一ページに「tertia」、「4」、「V」という番号が並んでいた。

アラビア数字とは直接関連がないものの、中世占星術的な数字符号についても触れなければならない。これらはいくつかの早期の著者によって記述されており、特にノヴィオマグス（1539年）において詳細に説明されている：

[図版]

次に、ラディカルの時代から十七世紀までのヨーロッパ全域でヒンゥー・アラビア数字がローマ数字を徐々に置き換えていく過程を見ることができます。しかし、現在の東アジアでは多くの国でこれらの数字はまだ知られておらず、その普及にはまだ道程があります。インドの一部、日本の庶民や中国の人々、シナム、マレーシア半島周辺、チベット、そして東インド諸島において、住民たちは依然として自国の数字を使用しています。しかしながら、西洋文明が東アジアの商業生活に進出し始めた段階で、これらの数字は徐々に定着してきました。ただし、インドの一部ではサンスクリット文字による表記も見られます。したがって、数学を学ぶ者にとっては驚くべきことに、これらの西欧で広く使用されている数字形式はいかに近代的であるか、現在でもその使用範囲はいかに限定的であるか、また世界がこのヒンダーラビア数字という単純な装置を受け入れるのにどれほど遅い道程をたどってきたかを理解する必要があります。

\clearpage

\end{document}
